% 本文件由 LaTeX 排版 Agent 根据有效 translation.json 自动生成。
% author=Miscellaneous; work_id=0835; title=伪教令集

\chapter{伪教令集}

% structure: p0001 source=h1 target=omitted reason=page-title-already-rendered-by-parent

% structure: p0002 source=h2 target=section reason=relative-heading-rank; base=h2

\section{则斐琳第一书函}

\Emph{致西西里全体主教。论主教审判及较重要教会案件之最终裁定在宗座。}\par

则斐琳，罗马城总主教，致所有居于西西里的诸位主教，在主内问安。\par

我们应当铭记天主对我们的恩宠，祂因自己的仁慈眷顾，为这目的将我们提升至司铎尊荣的顶峰，使我们遵守祂的诫命，并受命管理祂的某些司铎，得以禁止非法之事，教导当行之事。黑夜不能熄灭天上星辰，照样，世界的不义也不能蒙蔽那持守圣经可靠支撑的信徒之心。因此，我们应当深思熟虑，留心圣经以及其中所载的天主诫命，好使我们能表明自己不是违犯者，而是天主法律的实行者。\par

如今，宗主教和首席主教在调查一位被控告主教的案件时，不应作出最终判决，除非得到宗徒的权威支持，发现当事人要么自己认罪，要么由可靠且经正常审查的证人证明有罪，这些证人的人数不应少于主指示为协助宗徒而拣选的门徒数目——即七十二人。至于那些应凭天主权威铲除的诋毁者，以及仇敌的参谋（\Emph{auctores inimicorum}），我们不允许他们参与主教的指控或作证反对他们；也不应因下属的指控而控告或定任何上级的罪。在可疑的案件中，不应作出决定性判决；任何审判若非按秩序进行，便无效。此外，不得缺席审判任何人，因为神律和人律都禁止这样做。控告这些人的原告也应毫无嫌疑，因为主愿意祂的支柱稳固，不被任何随心所欲之人动摇。若非由他们的合法审判官作出判决，则不应约束他们中的任何人，因为连世俗法律也如此规定。因为任何被控告的主教，若有必要，可以自行挑选十二位审判官，由其公正审理他的案件。在挑选这些人之前，不应听取他的陈述、将他逐出教会或审判他；在他被正规传唤至由他本省主教组成的首次会议后，他的案件应由他们公正审理，并依据正确原则调查。然而，其案件的最终裁定应提交宗座，以便在那里作最后决定。正如古时由宗徒或其继承人所规定的那样，未经其权威认可，案件不得了结。所有案件，尤其是受压迫者，都应向宗座上诉并投靠它，如同投靠母亲，以便得到她的哺乳，受其权威保护，从压迫中得解脱，因为\BibleQuote{母亲不能，也不应忘记她的儿子。}\BibleRef{依 49:15} 因为主教审判和较重要教会案件，正如宗徒及其圣洁继承人所规定，应由宗座与其他主教共同最终裁定，而非由其他任何机构；因为尽管这些案件可移交给其他主教，但以下话却是对蒙福的宗徒伯多禄说的：\BibleQuote{凡你在地上所束缚的，在天上也要被束缚；凡你在地上所释放的，在天上也要被释放。}\BibleRef{玛 16:19} 此外，其他仅授予这圣座的特权，在宗徒及其继承人的宪章以及许多与此相符的其他文件中均可找到。因为宗徒连同许多其他主教制订了七十项法令，并规定予以遵守。因为鲁莽判断他人内心的隐秘是罪；因怀疑而责备那些行为看来并非不善的人，是不公义的，因为唯有天主才能判断人所不知的事。但祂\Quote{知道人心的隐秘}，而非他人。因为不义的审判是所有人都应提防的，尤其是天主的仆人。\BibleQuote{主的仆人不应争辩，}\BibleRef{弟后 2:24} 也不应伤害任何人。因为主教应由平信徒和神职人员忍受，主人应由仆人忍受，以便在忍耐的操练下，维持暂时之事，盼望永恒之事。因为不违反宗教宗旨的忍耐能增添美德的价值。你们应热切留意，不让任何弟兄受到严重伤害或毁灭。因此，你们应扶助受压迫者，从迫害者手中解救他们，以便能同蒙福的约伯一起说：\BibleQuote{那将亡者的祝福必临于我身，我使寡妇的心欣慰。我披上正义，以正义为袍和冠冕。我是盲人的眼，跛子的脚。我是穷人的父亲，我所不知的事，我仔细查考。我打碎了恶人的牙床，从他牙齿中夺出掠物。}等等。因此，你们这些被天主置于高位的人，应全力制止并击退那些为弟兄设陷阱、或对他们制造纷争和绊脚石的人。因为言语容易欺骗人，却不能欺骗天主。所以，你们应远离并提防这些人，直到这种黑暗完全消散，晨星照耀他们，喜乐升起，至圣的弟兄们。颁发于九月二十日，在至显赫的撒图尔尼诺和加利卡诺执政期间。\par

% structure: p0007 source=h2 target=section reason=relative-heading-rank; base=h2

\section{则斐琳第二书函}

\Emph{致埃及省主教。}\par

则斐琳，罗马城总主教，致在埃及事奉主的最亲爱的弟兄们。\par

我们从主——这神圣宗座和宗徒教会的建立者，以及从真福宗徒之长伯多禄——领受了如此大的信任，以致我们能以不懈的爱德为普世教会（已为基督的血所救赎）而劳苦，并以宗徒权威援助所有事奉主的人，帮助所有虔敬生活的人。凡在基督内（\BibleRef{弟后 2:24}）虔敬生活的人，必会遭受不虔敬者和外人的辱骂，被视作愚人和疯子，好使那些丧失现世美好事物以得永恒者，变得更美好、更纯洁。但那些折磨和嘲笑他们之人的轻蔑和讥讽，将回落在他们自己身上，当他们的富足变为匮乏，他们的骄傲变为羞辱。\par

% structure: p0011 source=h3 target=subsection reason=relative-heading-rank; base=h2

\subsection{一、论某些主教的劫夺或驱逐}

据你们的代表在宗座报告，我们的一些弟兄，即主教们，正被逐出他们的教会和座位，并被剥夺财产，且这样一无所有、被劫掠之后被传讯受审；此事完全无理，因为宗徒及其继任者的章程、皇帝的法令以及法律的条例都禁止这样做，宗座的权威也禁止这样做。事实上，在古代法令中已经规定，被驱逐并被剥夺财产的主教应恢复他们的教会，并且首先应全部归还他们的财产；其次，如果有人愿意公正地控告他们，他应承担同样的风险；法官应谨慎，主教应明智、和谐地共融于教会，在那里他们应为每一个看似受压迫的人作证；并且，在他们所有的一切以及他们的教会依法无损地归还之前，他们不应答辩。弟兄们，如果他们迫害你们，这并不奇怪，因为他们甚至曾迫害你们的元首、我们的主基督至死。然而，连迫害也应当忍耐，好使你们被认出是祂的门徒，你们也正是为祂而受苦。因此，祂亲自说：\Quote{为义受迫害的人是有福的。} \BibleRef{玛 5:10} 受这些见证的扶持，我们不应过分畏惧人的辱骂，也不应被他们的谴责击倒，因为主藉依撒意亚先知给我们这个命令说：\Quote{你们认识正义、心中有我法律的人民，请听我说：不要怕人的辱骂，也不要因他们的毁谤而惊惶；} \BibleRef{依 51:7} 并思想诗篇上所记载的：\Quote{难道天主不察知此事？因为祂洞悉人心的隐秘，知道这些人的思想不过是虚幻，} \Quote{他们彼此言谈虚妄，口吐谎言，内心诡诈。但上主必割除一切诈言，一切口出狂言的舌；他们曾说：我们的唇舌属于我们，谁为我们作主？} 因为如果他们记住这些事，就决不会爆发如此大的邪恶。他们这样做并非出于\OriginalTerm{可称赞的慈父般的教导}{probabili et paterna doctrina}，而是为了对天主的仆人发泄他们的报复情绪。因为经上记载：\Quote{愚人的道路，在他眼中看来正直；} \BibleRef{箴 12:15} 以及：\Quote{有一条路，人以为正直，但终途却是死亡之路。} \BibleRef{箴 14:12} 现在我们这些受这些苦的人，应当把这些事交托给天主的审判，祂要按各人的行为报应各人（\BibleRef{玛 16:27}）；祂也曾通过祂的仆人发出雷声说：\Quote{报复是我的，我必偿还。} \BibleRef{罗 12:19} 因此，你们要以信实彼此协助，并以行动和诚心互相帮助；谁也不要从帮助弟兄的事上收回手，因为主说：\Quote{如果你们彼此相爱，众人因此就可认出你们是我的门徒。} \BibleRef{若 13:35} 因此，祂也通过先知说：\Quote{看，弟兄们同居共处，多么美好，多么愉快！} 我解释这指的是属灵的居所，是在天主内的和谐，是在信仰的合一中，这信仰按真理区分这美好的居所，这确实在亚郎和那些穿上尊荣的司祭身上更美地体现出来，如同油膏在头上，滋养至高的理解力，并引领甚至达到智慧的终点。因为在这居所中，主已应许祝福和永生。因此，领悟到先知这话的重要，我们出于爱德说了这现在的兄弟之言，绝不寻求或意图寻求自己的利益。因为以毁谤还报毁谤，或\Emph{照俗语所说}\OriginalTerm{以桩击桩}{excutere palum palo}，是不好的。愿这事远离我们。这样的行径不属于我们。愿天主禁止。在天主公正的审判下，有时权力被赐给罪人迫害祂的圣徒，好使那些受天主之神扶助和承载的人，通过痛苦的锻炼而变得更加光荣。但那些迫害、辱骂、伤害他们的人，无疑将有祸。祸哉，祸哉，那些伤害天主仆人的人！因为对他们的伤害，关系到他们所事奉、所履行其职能的那一位。但我们祈祷，愿一道关闭的门放在他们嘴上，因为我们不愿任何人因他们的口而丧亡或被玷污，也不愿他们用口思想或散布任何有害的话语。因此，主也通过先知说：\Quote{我说：我要谨守我的道路，免得我以口舌犯罪。} 愿全能的主、祂的独生子、我们的救主耶稣基督，赐给你们这激励，使你们尽一切可能援助所有处于任何患难中的弟兄，并照所应当的，视他们的痛苦如己。用言语和行为给他们最大的帮助，好使你们被找到是祂的真门徒，祂曾命令众人彼此相爱如同爱自己。\par

% structure: p0013 source=h3 target=subsection reason=relative-heading-rank; base=h2

\subsection{二、论司铎与执事的圣秩}

此外，应在适当之时、众多证人面前隆重举行长老与肋未人的祝圣礼；应擢升经过考验、有学识之人担任此职，使你因他们的共融与帮助而大感喜悦。要不断将你内心的信赖置于天主的良善之上，并将这些以及其它神圣之言传于后世：\Quote{因为这就是我们的天主，直到永远，祂将引领我们直到永恒。}发于十一月七日，最尊贵的撒图尔尼诺与加利加诺执政之年。\par

% structure: p0015 source=h2 target=section reason=relative-heading-rank; base=h2

\section{卡利克斯特教宗第一封书信}

\Emph{致本笃主教。论四季斋期，及不得对导师提出控告。}\par

\Person{卡利克斯特}，罗马城天主教会总主教，致我们可敬的弟兄与主教\Person{本笃}，在主内问安。\par

因兄弟之爱的约束，并因宗座规训的催促，我们理当按主所赐予的，答复弟兄们的询问，并以宗徒封印的权威供给他们。\par

% structure: p0019 source=h3 target=subsection reason=relative-heading-rank; base=h2

\subsection{一、论斋期。}

你们素来知晓每年三次的斋戒，如今我们规定更为适宜的方式，于四季举行，即随着一年四季的轮转，我们也在四季中每季守一次隆重斋戒。正如我们因五谷、酒和油而身体得饱饫，同样，我们也当因守斋而灵魂得饱饫，符合先知匝加利亚的话，他说：\Quote{上主的话传给我说：万军的上主这样说：我决意惩罚你们时，因为你们的祖先激怒了我，我并未后悔；同样，在这些日子里，我决意善待耶路撒冷和犹大家，你们不要害怕。你们应做的事是：每人应对近人说实话，在你们的城门内应审判真实与和平的判决，不可心中图谋恶事陷害邻人，不可喜爱虚假的誓言，因为这一切都是我所憎恨的——万军的上主说。万军的上主的话又传给我说：万军的上主这样说：四月斋、五月斋、七月斋和十月斋，必要成为犹大家的喜乐、欢欣和愉快的节期；但你们应喜爱真理与和平——万军的上主说。}\BibleRef{匝 8:1}因此，我们应同心合意，依照宗徒的教导，众口一词，在我们中间没有分裂。让我们在同样心意、同样判断上达至成全；\EditorialNote{伯前三章}你以热忱成为这项工作的伙伴，我们为此欣喜。因为肢体与头不和是不相宜的；按圣经的见证，\EditorialNote{格前十二章}所有肢体都应跟随头。此外，无人怀疑宗徒教会是众教会之母，你们不应从她的规诫有丝毫偏离。正如天主子来是为承行父的旨意，同样，你们也应承行你们母亲——教会——的旨意，她的头，如前所述，是罗马教会。因此，凡违反此教会纪律、未经公正裁决所作之事，决不允许被视为有效。\par

% structure: p0021 source=h3 target=subsection reason=relative-heading-rank; base=h2

\subsection{二、论对导师的控告。}

此外，不得有人对一位教师（\OriginalTerm{教师}{teacher}）提出控告，因为儿子指责父亲、奴隶伤害主人是不对的。凡他们所教导的人，就是教师的儿子；如同儿子应爱其肉身之父，他们也应爱其灵性父亲。因为谁若信仰不正确，或指责父亲，或诽谤他们，他的生活就不正确。因此，教师（他们也被称为父亲）更应被容忍，而非被指责，除非他们背离了真信仰。因此，不得有人用\OriginalTerm{书面}{per scripta}控告教师；他也不得向任何控告者答辩，除非那人值得信赖、依法认可，且生活与言谈也无可指责。因为教师答复愚昧无知、生活可责之人，乃是相称于其愚昧的不配之事；正如圣经所说：\BibleQuote{不要按愚昧人的愚昧回答他。}\BibleRef{箴 26:4}谁若信仰不正确，他的生活就不正确。有信德的人不存恶意。若有人是\OriginalTerm{有信德的人}{a believer}，就当注意不捏造虚假指控，也不为任何人设下圈套。有信德的人常行事于信德；无信德的人则诡计多端，图谋使那些有信德、生活于虔诚与正义的人败坏，因为同类相求。无信德的人是在活着的身体中死去的人。另一方面，有信德之人的言论保护听众的生命。因为正如公教教师，尤其主的司铎，不应涉及任何错谬，同样他也不应受任何阴谋或激情的冤枉。圣经确实说：\BibleQuote{不要顺从你的欲情，要管束你的欲望；}\BibleRef{德 18:30}我们必须抵抗此世的许多诱惑和许多虚浮，以获得真正节制的完整，其首要瑕疵是骄傲，它是过犯的开端和罪的根源；因为怀有欲念的心志既不知节制，也不知献身于虔诚。善人没有敌人，除非是恶人，他们被允许如此，只是为了善人可以借此受纠正或受锻炼。因此，凡无可指摘的，皆为公教会所维护。无论君王，还是任何遵守虔诚的人，都不得胆敢做出违逆神圣诫命的事。因此，由法官因惧怕或奉君王、主教或有势力者的命令而作出或执行的任何不义判决或不义\OriginalTerm{裁定}{diffinitio}，均属无效。虔诚人不应当仅仅满足于不介入他人的仇怨，或不以恶言增加之，除非他也致力于以善言息灭仇怨。在恶行中谦卑告解，胜过在善行中骄傲自夸。此外，所有过蒙福生活的人，宁愿在和平与正义的合宜状态下走这条道路，也不愿使自己陷入我们罪恶的报应之苦。因为我记得，我是在那位的名下治理教会，他的宣认曾由我们的主耶稣基督所尊崇，他的信德常摧毁一切错谬。而且我知道，我只能将全部精力用于此事业，因为普世教会的福祉在此\OriginalTerm{危在旦夕}{infestatur}。我也希望天主的仁慈如此眷顾我们，以便藉着祂的宽仁，每一致命的疾病得以除去——由天主亲自驱除——并在祂的启示与助佑下，凡有益的事得以成就，以赞美你们的信德与热忱。因为若非司铎的权威在神圣职务的服务范围内支持它们，一切事物都无法安全。颁于十一月二十一日，最杰出的\Person{安多尼诺}与\Person{亚历山大}担任执政官期间。\par

% structure: p0023 source=h2 target=section reason=relative-heading-rank; base=h2

\section{教宗卡利斯督第二封书信}

\Emph{致高卢全体主教。论不可参与阴谋及其他非法行径，以及忏悔后失足者的恢复。}\par

\Person{卡利斯督}致我们最亲爱的弟兄，即安顿在高卢各地的全体主教。\par

根据许多人的报告，我们得知，你们的仁爱藉着圣神的热情，面对一切攻击，如此坚定地掌握并引领教会的船舵，以致在天主的旨意下，它既未遭遇船难，也未遭受船难的损失。因此，我们为这些见证而欣喜，恳请你们不要让那些地区做出任何违反宗徒规条的事；但要依靠我们的权威，制止有害之事，禁止非法之事。\par

% structure: p0027 source=h3 target=subsection reason=relative-heading-rank; base=h2

\subsection{一、论那些阴谋反对主教或与这些人同谋的人。}

现我们听说，在你们那地方盛行阴谋的罪行，且据我们得知，人民正合谋反对他们的主教；此罪的伎俩不仅在基督徒中令人憎恶，就连在异教徒中也如此，且为外邦法律所禁止。因此，不仅教会的法律，甚至世俗的法律都谴责犯此罪的人；不但那些实际密谋者，就连那些参与其中者，也受谴责。\BibleRef{罗 1:32}此外，我们的前任们会同为数众多的主教规定，凡处于司铎尊位并属于圣职人员中\Emph{犯此罪者}应被剥夺他们所享有的尊荣；并命令其他人应被切断共融，逐出教会；他们同时下令，两种等级中所有人都应被宣布为\OriginalTerm{可耻的}{infames}；而且不仅对行此事者，也对那些参与其中者。因为，轻视天主的命令、证明自己不顺从教父训令的人，理应受到更严厉的惩罚，好使他人畏惧行此类事，并使众人皆在兄弟般的和睦中喜乐，且都为严厉与良善的榜样所感召。因为，若我们（但愿天主禁止）忽视了教会的关怀，不顾其力量，我们的懈怠将摧毁纪律，并必然损害信徒的灵魂。此外，这种人不得被准许控告任何人：无论是他们的声音，还是那些受限制者的声音，都不能伤害或指控任何人。\par

% structure: p0029 source=h3 target=subsection reason=relative-heading-rank; base=h2

\subsection{二、论与受绝罚者或不信者交往的人）}

同样，对于那些受司铎绝罚的人，在双方公正审查之前，任何人都不得接纳他们；也不得与他们有言语、饮食、口亲致敬或任何形式的交往，也不得向他们致候；因为，凡明知故犯，在上述或其他禁止之事上与受绝罚者交往的人，按照宗徒的规定，将使自己受到同样的绝罚。因此，圣职人员和教友应当远离这些人，以免承受相同的惩罚。也不应与不信者联合，不要与他们有任何共融。事实上，做这些事的人不被视为信徒，而是被视为不信者。因此宗徒说道：\BibleQuote{信者与不信者有什么相干？义与不义有什么相通？}\BibleRef{格后 6:14}\par

% structure: p0031 source=h3 target=subsection reason=relative-heading-rank; base=h2

\subsection{三、没有主教应干预属于其他教区的事，以及论主教的调动）}

再者，任何人不得侵犯他人的界限，也不得擅自审判或绝罚属于他人堂区之人；因为这样的审判、任命、绝罚或定罪，既不得被认可，也毫无效力；因为任何人不应受其本身审判者以外之人的判决所约束，也不应被其定罪。因此主亦言及此事：\BibleQuote{不可挪移你祖先所立的古界。}\BibleRef{箴 22:28} 再者，任何首席主教或都主教不得\Emph{侵入}教区（主教）\OriginalTerm{教区（主教）}{diœcesani}的教会或堂区，也不得擅自绝罚或审判属于其堂区的任何人，或未经其商议或判决而行事；但他应遵守此项法律，即由宗徒、教父、我们的前任所制定，并由我们批准的法律：即任何都主教，除仅限于其本堂区之事外，若未经所有同省主教的商议和同意而试图行事，他将冒失去其地位的风险，且如此行事应视为无效；但就省主教团体的事务及其教会、圣职人员和平信徒的需要所必须安排或处理之事，应经同省所有主教同意而行，且不可有丝毫的统治傲慢，而应以最谦卑和谐的行动，正如主所说：\BibleQuote{我来不是为受服侍，而是为服侍。}\BibleRef{玛 20:28} 在另一处祂说：\BibleQuote{你们中谁愿为大，就当作你们的仆役，}\BibleRef{谷 10:44} 等等。同样，同省的主教们也应在一切事上与之商议而行，除非涉及他们本堂区之事，应遵照圣教父们的规章（他们虽比我们先与世长辞，却从同一纯洁泉源汲取真理与信德之光，藉同一启迪和引导的圣神，寻求天主教会之繁荣及所有基督徒的共同益处），以便一心一口、同心合意，使天主圣三永受光荣。任何首席主教、都主教或其他主教，均不得擅自进入他人的教座，或占据不属于他、而属于另一位主教堂区一部分的产业，无论受何人指示，除非他受到公认管辖权所属之人的邀请；若他愿保持其地位之尊荣，则不得在此进行任何安排、指令或审判。但若他擅自妄为，将被定罪；不仅他本人，连同合作及赞同者亦然。因为正如在此等情况下\OriginalTerm{任命}{ordinatio}之权力被禁止，同样，审判或处理其他事务之权力亦被禁止。因为若人无权任命，又岂能审判？无疑，他绝不能审判，也无权审判：因为正如他人之妻不得与任何人\OriginalTerm{通奸}{adulterari}，且除其丈夫生前外，任何人均不得审判或处置她；同样，绝不允许在主教（其妻子无疑指其教会或堂区）生前，未经该主教的审判和善意，由他人审判或处置其妻，或享受他人之怀抱，即受他人祝圣。因此宗徒说：\BibleQuote{丈夫活着的时候，妻子是受法律约束的；如果丈夫死了，她就从丈夫的法律中解脱了。}\BibleRef{罗 7:2} 同样，主教的配偶（因为教会被称为他的配偶和妻子）在他生前受其约束；他死后，她便得解脱，可随意嫁人，但只可在主内，即按规矩而行。因为若她丈夫活着时嫁与他人，她将被视为通奸者。同样，他若自愿另娶，也被视为通奸者，并被剥夺共融的特恩。然而，若他在自己的教会受迫害，他必须逃往另一教会，并与她结合，正如主说：\BibleQuote{如果他们在这城里迫害你们，就逃到别的城里去。}\BibleRef{玛 10:23} 然而，若变更是为了教会的\Emph{益处}，他不可自行其事，而须应弟兄们的邀请，并得此圣座之准许，且不可出于野心，而应为公众利益。\par

% structure: p0033 source=h3 target=subsection reason=relative-heading-rank; base=h2

\subsection{四、论血亲婚姻及由之所生之人，以及法律所驳回的控告。}

此外，血亲婚姻是被禁止的，因为一切法律，无论神圣还是世俗，都禁止此类。因此，神圣法律不仅驱逐，甚至绝罚那些如此行的人以及他们所生的后代。世俗法律则称此类人为声名狼藉者，并禁止他们继承遗产。我们也追随教父，紧随其足迹，以恶名标记他们，并视其为声名狼藉者，因为他们沾染了恶名的污点。我们也不应接纳那些世俗法律所驳回的人或其控告。（因为谁怀疑，若人法律与理性和荣誉不矛盾，便应予采纳，尤其当它们促进公共利益或维护教会职务的权威，并予以支持时？）我们所称的血亲，是指神圣法律及罗马和希腊皇帝的法律称为血亲的人，并且承认其继承权，不能将其排除。因此，此类婚姻既不合法，也不能有效，而应予以拒绝。（若有人胆敢尝试，则应由宗座权威予以撤销。）\par

% structure: p0035 source=h3 target=subsection reason=relative-heading-rank; base=h2

\subsection{五、论那些不应被接纳提出控告或作证之人；以及证据仅应就当事人面前发生之事提出。}

因此，凡未经合法结婚，或没有聘礼头衔（\OriginalTerm{聘礼头衔}{dotali titulo}）与司铎祝福而结合的人，决不能控告司铎或那些合法结婚的人，也不能对他们作证，因为凡沾染乱伦污点的人都是可耻的，不得控告上述之人。因此，不但他们，连所有同意他们的人，也应被排斥，并被视为可耻。我们认为，对强盗或殴打长者的人也应如此。世俗的法律固然处死这种人；但对我们来说，仁慈居首位，我们接受他们带着可耻的标记去悔改。我们无法除去他们所沾染的耻辱；但我们的愿望是通过公开补赎和向教会做出补偿来医治他们的灵魂：因为公开的罪过不能通过秘密的纠正来洗净。此外，凡在正确信仰上可疑的人，决不应被允许控告司铎和那些信仰无疑的人；在人类作证的事情上，这样的人应被视为权威可疑。因此，信仰可疑之人的声音应被视为无效；对不认识正确信仰的人，不应给予信任。因此，在审判中，应查问控告者和被控告者的品行和信仰；因为品行和信仰不正确、其生活可受指控的人，不得控告他们的长辈；对那些信仰、生活和自由不明的人，也不能给予这种许可。也不应准许卑鄙的人控告他们。但要仔细查究控告者的人品（\OriginalTerm{要仔细查究控告者的人品}{rimandae sunt enucleatim personae accusatorum}）；因为不经书面文件不应轻易接纳他们，且绝不可仅凭书面文件接纳他们作为控告者。因为没有人可以仅凭书面文件控告或被控告，而必须用活的声音；每个人都必须在他要控告的人面前提出控告。不在被控告者面前时，不应相信任何控告者。同样，证人也不得仅凭书面文件作证；他们必须亲自忠实地作证，且只对他们亲眼看见并知道的事情作证。他们不得在任何其他案件或事情上作证，只限于已知发生在他们面前的事情。此外，同一血统的控告者不得对非亲属作证，家仆（\OriginalTerm{家仆}{familiares}）或来自同一家庭的人也不得；但如果他们愿意且彼此同意，只有父母应在这些案件中作证，其他人不得。凡有任何可疑之处的控告者或证人都不得被接纳；因为亲属、友谊或主仆关系往往会阻碍真理。肉体的爱、恐惧和贪婪通常会使人的感知迟钝，歪曲他们的判断；以致他们将利益视为敬虔，将金钱视为谨慎的报偿。因此，任何人都不可用诡诈对邻舍说话。恶人的口是深坑。无辜的人轻易相信，就容易跌倒；但即使跌倒，他还会起来；而诡诈的人用尽伎俩却直坠毁灭，永不能起来或逃脱。因此，每个人都当谨慎言语，不可对别人说对自己不愿说的话。因此，圣经说得好：\BibleQuote{你不愿意别人对你做的，也不要对别人做。} \BibleRef{多 4:15} 因为我们做任何完美之事都需要时间（\OriginalTerm{更为成熟}{maturius}）；我们不可在计划或工作中急躁，也不可违反秩序。但若有人在某事上跌倒，我们不可将他推入毁灭；而要以兄弟之爱责备他，正如蒙福的宗徒所说：\BibleQuote{若有人偶然被过犯所胜，你们属神的人要用温柔的心把他挽回；又当自己小心，恐怕也被引诱。你们各人的重担要互相担当，如此就完全了基督的法律。} 此外，圣洁的\Person{达味}曾犯下致命的罪而需悔改，但仍保留了尊荣。蒙福的\Person{伯多禄}在否认主后痛哭悔改，但仍是宗徒。主也通过先知向犯罪者应许：\BibleQuote{恶人若回头离开所行的恶，行正义与公道，我必不再记念他所犯的一切过犯。} \BibleRef{则 18:21}\par

% structure: p0037 source=h3 target=subsection reason=relative-heading-rank; base=h2

\subsection{六、论司铎在失足后是否仍可施行圣事}

因为那些认为主的司铎在跌倒之后，即使他们可能表现出真正的悔改，也不能再侍奉主、从事他们尊贵的职务，即使此后他们能度圣善的生活并正确持守司铎职的人，是陷于错误的。那些持此意见者不仅陷于错误，而且似乎也在争议并行事反对托付给教会的\Emph{钥匙的权能}，关于此权能经上记载说：\BibleQuote{凡你们在地上所释放的，在天上也要被释放。} \BibleRef{玛 18:18} 简而言之，这见解要么不是主的，要么就是真的。但尽管如此，我们毫不迟疑地相信，主的司铎和其他信徒，在为他们的错误作了适当的补赎之后，都可以恢复他们的荣耀，正如主自己藉祂的先知所证实的：\BibleQuote{那跌倒的岂能不再起来？那转离的岂能不再回来？} \BibleRef{耶 8:4} 在另一处，主说：\Quote{我不喜悦罪人死，而喜欢他归正并生活。} 先知达味在悔改时说：\Quote{求祢将祢救恩的喜乐还给我，并以祢慷慨的圣神扶持我。} 他确实在悔改后也教导别人，并向天主奉献祭献，从而给圣教会的教师树立了榜样：如果他们跌倒，其后向天主作出了正确的悔改，他们也可以同样地做这两件事。因为他教训时曾说：\Quote{我要将祢的道路指教悖逆之辈，罪人必归向祢。} 他为自己献祭时也说：\Quote{天主所要的祭献就是忧伤痛悔的心。} 因为先知看见自己的过犯藉着悔改得洁净，就毫不怀疑藉着宣讲和向天主献祭也能医治他人的过犯。因此，流泪能触动心灵的\OriginalTerm{情感}{passionem}。当补赎圆满完成时，心就会转离愤怒。因为那人不宽恕自己的邻人，又如何想自己能获得仁慈呢？如果过犯增多，就让仁慈也增多罢；因为上主有仁慈，祂有丰富的救赎。主的手中有一切的丰盛，因为祂是\OriginalTerm{权能}{virtutum}的主，荣耀的君王。因为宗徒说：\BibleQuote{众人都犯了罪，缺乏天主的荣耀；但因天主的恩宠，藉着在耶稣基督内的救赎，白白地成义。天主立了祂为赎罪祭，藉信德在祂的血中，为显示祂的正义，因为祂暂且容忍了先前的罪恶；为显示祂的正义，祂在现今的时期，表示祂是正义的，又使那信耶稣的人成义。} \BibleRef{罗 3:23} 达味说：\Quote{罪过得赦免、罪恶得遮盖的人，是有福的。} 因此，人虽已跌倒，却因天主的恩宠得以洗净罪污，重新起来，并依照上述权威的见证，留在原先的地位。他应当小心不再犯罪，使福音的裁断常在他身上：\BibleQuote{去吧，不要再犯罪了。} \BibleRef{若 8:11} 因此宗徒说：\BibleQuote{所以，不要让罪在你们必死的身体上为王，致令你们顺从它的情欲；也不要把你们的肢体交给罪，作为不义的武器；但要把你们自己交给天主，犹如从死者中复活的人，并把你们的肢体交给天主，作为义德的武器。因为罪不应再统治你们，因为你们已不在法律之下，而在恩宠之下。那么，我们因为不在法律之下，而在恩宠之下，就可以犯罪吗？绝对不可！难道你们不知道：你们将自己献给谁当奴仆，服从他，就成了他的奴仆？或作罪恶的奴仆，以致死亡；或作顺命的奴仆，以致成义。感谢天主，你们从前虽然作了罪恶的奴仆，现今却从心里服从那传授给你们教理的标准。你们既脱离了罪恶，就成了义德的奴仆。我按人的常情来说话。} \BibleRef{罗 6:12} 因为判断人的罪比被判断的罪更大。\BibleQuote{你以为，}宗徒说，\BibleQuote{人啊！你判断那些作这些事的人，自己也作同样的事，你以为你能逃脱天主的审判吗？或者你轻视祂丰厚的慈爱、宽容和忍耐，不晓得天主的慈爱是领你悔改吗？你因着顽梗不悔的心，就在你发怒和显示公义审判的日子，为自己积蓄忿怒。祂要按各人的行为报应各人：凡坚持行善，追求光荣、尊贵和不朽的，赐以永生；凡固执作恶，不服从真理，而服从不义的，报以忿怒和愤恨，患难和困苦，加于一切作恶的人，首先为犹太人，也为希腊人；但光荣、尊贵与平安，赐给一切行善的人。} \BibleRef{罗 3:3} 我的弟兄们，你们不仅应避免持有，甚至应避免听取那排除仁慈的判决；因为仁慈胜过一切全燔祭和牺牲。 \BibleRef{谷 12:33} 我们简短地答复了你们的问题，因为你们的信札来时我们负担过重，且忙于其他判决。于十月八日，在最为显赫的\Person{安多尼诺}与\Person{亚历山大}的执政时期颁发。\par

% structure: p0039 source=h2 target=section reason=relative-heading-rank; base=h2

\section{教宗乌尔巴诺一世的书信}

\Emph{致全体基督徒。论及教会只接受信徒的财产，而不接受其价值，如同宗徒时代一样；以及为何应在教堂内为主教预备高座；以及关于主教所开除教籍的人，任何人都不得与之交往，且无论以何种方式，均不得接纳被主教驱逐的人。}\par

乌尔巴诺主教，祝圣于圣神，顺服于主耶稣基督的宝血，致全体基督徒，问候。\par

最亲爱的诸位，所有基督徒都应当效法他们领受其名的那一位。宗徒雅各伯说：\BibleQuote{我的弟兄们，若有人说自己有信德，却没有行为，有什么益处？}\BibleRef{雅 2:14} \BibleQuote{我的弟兄们，不要多数人当师傅，因为你们知道，\OriginalTerm{你们领受}{sumitis}更严厉的审判；因为众人都犯许多过失。}\BibleRef{雅 3:1} \BibleQuote{你们中谁是有智慧，有见识的呢？让他以良好的品行，和智慧的温和，显示他的行为。}\BibleRef{雅 3:13}\par

% structure: p0043 source=h3 target=subsection reason=relative-heading-rank; base=h2

\subsection{一、论共同生活，以及教会开始拥有财产的原因。}

我们知道你们并非不知道，迄今为止，在善良的基督徒中，一切财物共用的原则一直有效地实行，并且因天主的恩宠至今仍然如此；尤其是在那些被选为主之份的人当中，即圣职人员，正如我们在宗徒大事录中读到的：\BibleQuote{众信徒都一心一意，凡各人所有的，没有人说是自己的，都归公用。宗徒们以大德能，为复活的主耶稣基督作证，在众人面前大受爱戴。在他们中，没有一个贫乏的人，因为凡有田产房屋的，都卖了，把卖得的价钱带来，放在宗徒们脚前，照每人所需要的分配。有位若瑟，宗徒称他为巴尔纳伯，即解说\InnerQuote{安慰之子}，是肋未人，生于塞浦路斯，他有田地，也卖了，把银钱带来，放在宗徒脚前。}\BibleRef{宗 4:32}等等。因此，当大司祭和其他人，以及肋未人和其余的信友察觉，如果他们将要出售的遗产和土地交给主教所管辖的教会，或许更有益处，因为这些财产的收入可以为现在和将来持有共同信仰的信友提供更大更好的维持，胜过一次性卖得的金钱，他们便开始将惯常出售的财产和土地交给母教会，并习惯于靠这些收入过活。\par

% structure: p0045 source=h3 target=subsection reason=relative-heading-rank; base=h2

\subsection{二、论管理教会财产的人及其用途，以及侵犯者。}

此外，各堂区所拥有的财产交由主教保管，主教身居宗徒之位；此事至今如此，且当永远如此。主教与信友作为其管家，应当尽其所能向所有愿意度共同生活的人供应一切必需品，使他们中间无人匮乏。因为信友的财产也称为献仪，因它们是献给主的。因此，这些财产不应用于教会用途之外，也不应用于上述基督徒弟兄及穷人的利益之外；因为它们是信友的献仪，是罪的赎价（\OriginalTerm{罪的赎价}{pretia peccatorum}），是穷人的产业，且以上述名义交给了主。若有人另行处理（愿主禁止），他当小心，免得遭受\Person{阿纳尼雅}和\Person{撒斐辣}的惩罚，并被定为渎圣罪，如同那些对所指定财产的价银说谎的人一样，我们在上述\WorkTitle{宗徒大事录}的段落中读到这样的事：\BibleQuote{但有一个名叫\Person{阿纳尼雅}的人，同他的妻子\Person{撒斐辣}卖了田产（\OriginalTerm{田地}{agrum}），把价银私自留下几分，他的妻子也知道，其余的几分拿来放在宗徒脚前。\Person{伯多禄}对他说：\InnerQuote{阿纳尼雅！为什么\Person{撒殚}充满了你的心（\OriginalTerm{试探}{tentavit}），使你欺骗圣神，把田地的价银私自留下几分呢？田地还没有卖，不是你自己的吗？既卖了，价银不是你作主吗？你怎么心里起这意念呢？你不是欺骗人，乃是欺骗天主。}阿纳尼雅听见这话，就仆倒断了气。凡听见的人都甚惧怕。有些少年人起来，把他包裹（\OriginalTerm{移开}{amoverunt}），抬出去埋葬了。约过了三小时，他的妻子进来，还不知道这事。\Person{伯多禄}对她说：\InnerQuote{你告诉我，你们卖田地的价银就是这些吗？}她说：\InnerQuote{就是这些。}\Person{伯多禄}说：\InnerQuote{你们为什么同心试探主的圣神呢？埋葬你丈夫之人的脚已到门口，他们也要把你抬出去。}妇人立刻仆倒在\Person{伯多禄}脚前，断了气。那些少年人进来，见她已经死了，就抬出去，埋在她丈夫旁边。全教会和听见这事的人都甚惧怕。} \BibleRef{宗 5:1} 弟兄们，这些事应当谨慎提防，大大惧怕。因为教会的财产，不像私人财产，而像公共财产，且是献给主的财产，应当以极深的敬畏、本着忠信的精神来管理，且只为上述用途，免得那些将它从所托付的手中挪作他用的人犯下渎圣罪，免得他们遭受\Person{阿纳尼雅}和\Person{撒斐辣}的惩罚与死亡，更（更糟的是）免得他们成为\OriginalTerm{被诅咒者}{anathema maranatha}，并且，虽然他们的身体不像\Person{阿纳尼雅}和\Person{撒斐辣}那样仆倒死亡，但他们的灵魂——比身体更高贵的灵魂——却会死亡，被从信友的团体中剪除，沉入深渊。因此，所有人都必须留意此事，忠心守望，避免这种篡夺的羞辱，免得那些被奉献用于隐秘（或神圣）和天上之事的财产被任何侵入者掠夺。若有人这样做，那么，在对此罪行施以应得的严厉惩罚——须公正地施于渎圣者——之后，他当被判处永久耻辱，投入监牢，或被终身流放。因为，按照宗徒的话（\BibleRef{格前 5:5}），我们应当把这样的人交与\Person{撒殚}，使他的灵魂在主的日子可以得救。\par

% structure: p0047 source=h3 target=subsection reason=relative-heading-rank; base=h2

\subsection{三、关于任何人企图剥夺教会持有财产的权利。}

因此，由于上述的增加和生活方式，主教所管辖的教会因主的帮助而大大增长，如今大部分教会拥有如此多的财产，以至于在他们中间，没有人因选择共同生活而陷于贫穷；相反，这样的人从主教及其辅理人员那里得到一切必需品。因此，若有人现在或将来起来企图挪用这些财产，他当遭受已经提到的审判。\par

% structure: p0049 source=h3 target=subsection reason=relative-heading-rank; base=h2

\subsection{四、关于主教的座位。}

此外，关于在主教教堂中设有高起的座位，如同宝座一样设置好，这些显示主教被主赋予了监督和审判的权力，以及释放与束缚的权柄。因此，救主自己在福音中说：\BibleQuote{凡你们在地上所束缚的，在天上也要被束缚；凡你们在地上所释放的，在天上也要被释放。}\BibleRef{玛 18:18} 另一处：\BibleQuote{你们领受圣神吧！你们赦免谁的罪，谁的罪就得赦免；你们保留谁的罪，谁的罪就被保留。}\BibleRef{若 20:22}\par

% structure: p0051 source=h3 target=subsection reason=relative-heading-rank; base=h2

\subsection{五、没有人应与主教没有交往的人交往，也不应接纳主教所拒绝的人。}

最亲爱的诸位，我们把这些事摆在你们面前，好使你们明白主教们的权柄，并在他们身上敬畏天主，爱他们如同自己的灵魂；并使你们不与那些他们不与交往的人交往，也不接纳他们所驱逐的人。因为主教的审判是极其可怕的，即使他可能不公正地束缚某人，然而他却应当极其谨慎地避免这种事。\par

% structure: p0053 source=h3 target=subsection reason=relative-heading-rank; base=h2

\subsection{六、关于圣洗圣事中的誓约，以及那些献身于共同生活的人。}

我们在劝勉你们的同时，也劝告所有接受了基督信德、并且从基督那里取得了基督徒之名的人，千万不要使你们的基督徒生活变得空虚，而要坚定不移地持守在圣洗中所接受的承诺，这样你们才不致被弃绝，反而在他面前堪当。如果你们中有人进入了凡物公有的生活，并且立誓不拥有私有财产，就当注意不要使自己的诺言落空，而要忠信地遵守这向主所做的承诺，这样他为自己获得的不是惩罚，而是赏报；因为人若根本不立誓，比立了誓却不能尽力履行更好。因为那些立了誓或接受了信德，却没有遵守誓言，或者以恶事度过一生的人，所受的惩罚比那些没有立誓而生活，或没有信德而行善死去的人更重。为此，我们藉着本性天赋的理性心智，以及第二次出生的更新，为的是按照宗徒的教导，我们能够辨别（\OriginalTerm{辨别}{sapiamus}）上面的事，而不是地上的事；\BibleRef{哥 3:2}因为此世的智慧在天主面前是愚妄。\BibleRef{格前 3:19}最亲爱的诸位，此世的智慧除了教我们追求有害的事物、喜爱将要灭亡的事物、忽视有益的事物、轻视恒久的事物，还能做什么呢？它称赞贪爱钱财，经上说：贪爱钱财是万恶的根源；\BibleRef{弟前 6:10}而且它有一个特别的恶处：它虽然突出了短暂之物，却遮蔽了永恒之物；它只关注外在，却不查看内在隐藏的；它追求陌生的事物，却成为行恶者自身的陌生之恶。看，此世的智慧教导人什么呢？教人生活在享乐中。经上因此说：那生活在享乐中的寡妇，虽活犹死。\BibleRef{弟前 5:6}它教人以最柔软的快乐、罪恶、恶习和情欲来喂养肉身，以饮食无度来压迫灵魂，抑制属灵的生命，把刀剑交到敌人手中攻击自己。看，此世的智慧所给的是什么劝告？它叫好人宁愿选择作恶，叫他们心志错谬，渴望犯罪，而且不思想天主那可畏的声音，那时恶人将像草一样被焚烧。\par

% structure: p0055 source=h3 target=subsection reason=relative-heading-rank; base=h2

\subsection{七、论主教的覆手}

因为所有信徒在领受圣洗后，都应当藉着主教的覆手领受圣神，这样他们才能成为完全的基督徒；因为当圣神倾注在他们身上时，信者的心会因明智和坚定而扩大。我们领受圣神，为能成为属神的人；因为属血气的人不接受天主圣神的事。\BibleRef{格前 2:14}我们领受圣神，为能明智地辨别善恶，爱正义，憎恶不义，从而抵挡骄傲和恶意，抵抗奢侈和各种诱惑，以及不洁和可耻的欲望。我们领受圣神，为能藉着对生命的爱和对光荣的热忱，将我们的心思从地上的事物提升到天上和神圣的事物。——颁于九月诺涅日，即同月第五日，在最显赫的安托宁和亚历山大执政之年。\par

% structure: p0057 source=h2 target=section reason=relative-heading-rank; base=h2

\section{教宗\Person{Pontianus}的第一封书信}

\Emph{致\Person{Felix Subscribonius}，论应给予司铎的尊荣。}\par

\Person{Pontianus} 主教，致 \Person{Felix Subscribonius} 问安。\par

我们的心因你们的良善而极其喜乐，因为你们竭尽全力实行神圣宗教的实践，并坚固忧伤困苦的弟兄于信德与信仰中。因此，我们恳求我们救赎主的仁慈，愿祂的恩宠在一切事上支持我们，并使我们能实际完成祂所赐予我们渴望的事。因此，在这件善事上，酬报的益处正按我们对工程的热情之增长而倍增。由于在这一切事上，我们需要天主恩宠的帮助，我们以不断的祈祷恳求全能天主的仁慈，愿祂既赐给我们这些善工的渴望——这些善工应常由我们实行——又赐给我们行善的力量，并引导我们走向行善的果实——这\Emph{道路}乃是牧者之牧者所自称的——好使你们能藉着祂（离了祂一无所能）完成你们已开始的善工。此外，关于主的司铎们，我们听说你们帮助他们对抗恶人的阴谋，并支持他们的事业，要知道你们这样做极悦乐天主，祂召叫他们服侍自己，并以如此亲密的共融尊崇他们，好使藉着他们，祂接受他人的祭品，赦免他们的罪过，并使他们与自己和好。他们也用自己的口制作主的身体（\OriginalTerm{用自己的口制作主的身体}{proprio ore corpus Domini conficiunt}），并将之给予百姓。因为论到他们，经上说：\Quote{伤害你们的，就是伤害我；凡亏待你们的，必因他所行的不义而受报。}别处又说：\Quote{听从你们的，就是听从我；轻视你们的，就是轻视我；轻视我的，就是轻视那派遣我来的。}\BibleRef{路 10:16}因此，他们不应受骚扰，而应受尊敬。在他们身上，主自己受到尊敬，因为他们执行主的使命。因此，如果他们偶然跌倒，应由信友们扶起并支持。再者，如果他们犯罪，应由其他司铎指控；此外，应由首席主教们约束（\OriginalTerm{约束}{constringantur}），而不得由世俗人或生活邪恶者指控或约束。其次，他们不应受声名狼藉者、邪恶者、仇敌、或另一教派或宗教之成员的控告。因此，我们听到你们因兄弟的离世（\OriginalTerm{离世}{transitu}）而悲伤，我们并非不忧伤。为此，我们恳求全能天主以祂恩宠的嘘气（\OriginalTerm{嘘气}{aspiratione}）安慰你们，并以天上的守护保护你们不受恶灵和邪恶人的侵扰。因为如果你们在兄弟离世后不得不承受某些敌对者的扰乱，不要觉得奇怪，虽然你们寻求在自己的家乡——即在活人之地——享福，却在异地遭受人的恶事。因为现世生活是一种旅居；凡叹息寻求真正故乡的人，旅居之地便是试炼，无论它看似多么愉快。至于你们寻求故乡，在你们发出的叹息中，我也听到人类压迫的呻吟升起。这事因全能天主奇妙的安排而发生，为使真理以爱召叫你们时，现世通过它所带的苦难将你们的感情从自身抛回，心灵便更容易摆脱对现世的爱，因为它在被召时也被推动。所以，正如你们已开始的，请留意款待之责；在祈祷与眼泪中极其迫切地劳苦；现在更大方更自由地献身于你们一向喜爱的施舍，以便在酬报时，你们工作的收益与你们在此处勤勉的等级成正比例增长。\Signature{—— 颁于二月朔日前第十日（一月二十三日），在最有声望的塞维鲁斯与昆提安努斯执政期间。}\par

此外，我们以慈父般的愉悦问候你们的良善，恳求你们不要懈怠已开始的善工。愿无人能使你们远离它们；反之，愿圣职人员和天主的仆人们，以及所有旅居那些地区的基督徒们，能因基督和圣伯多禄的爱，充分发现你们爱德在一切事上的倾向，并在可能发生的每项需要中获得你们眷顾的安慰；以便所有人都能得到你们的援助的保护和帮助，我们也感激你们，我们的主耶稣基督将赐予\Emph{你们}永恒的荣耀，蒙福的宗徒之长伯多禄，你们为他而劳苦，愿他开启同一荣耀之门。\Signature{—— 颁于二月朔日前第十日（一月二十三日），在最有声望的塞维鲁斯与昆提安努斯执政期间。}\par

% structure: p0062 source=h2 target=section reason=relative-heading-rank; base=h2

\section{教宗\Person{彭西安}的第二封书信}

\Emph{给所有主教。论弟兄友爱与避恶。}\par

\Person{彭西安}，圣而公教会的主教，致所有正确敬拜主并热爱神圣敬礼的人，祝你们平安。\par

天主受享光荣于高天，主爱的人在世享平安。\BibleRef{路 1:14}最亲爱的，这些话不是人的话，而是天使的话；它们不是由人的理智所设计，而是在救主诞生时由天使所宣告。从这些话，所有人都能无疑地明白：主赐予平安，不是给心怀恶意的人，而是给善意的人。因此，主藉先知说：\Quote{天主对待心地纯洁的人何其美善！但至于我，我的脚几乎失稳，我的脚步几乎滑倒，因为我嫉妒恶人，见罪人享平安。}然而，真理亲自论及善人说：\BibleQuote{心里洁净的人是有福的，因为他们要看见天主。}\BibleRef{玛 5:8}那些思想邪恶或有害弟兄的事的人，并不是心里洁净的人；因为忠诚的人不图谋任何邪恶。因此，忠诚的人宁愿听适当的话，而不说不适当的话。若有人是忠诚的，就应谨守自己不说恶言，不为任何人设下陷阱。这样，天主的子女与魔鬼的子女便得以区分。因为天主的子女常思想并努力做属于天主的事，不断帮助弟兄，不愿伤害任何人。反之，魔鬼的子女则常图谋邪恶与有害的事，因为他们的行为是恶的。论到他们，主藉耶肋米亚先知说：\BibleQuote{我要向他们宣布我的判决，针对他们的一切邪恶。}\BibleRef{耶 1:16}\BibleQuote{因此，我还要与你们辩论——上主说——并与你们的子子孙孙辩论。}\BibleRef{耶 2:9}\BibleQuote{看，我计划降灾给你们，图谋打击你们。}\BibleRef{耶 18:11}弟兄们，这些事是极可畏惧的，所有人都应提防；因为天主的审判临到的人，不能安然离去。因此，每个人都应谨慎注意：他不可图谋或加害弟兄他自己不愿承受的事。有信仰的人甚至不可被怀疑说了或做了他自己不愿忍受的事。因此，可疑之人、敌意之人、好讼之人、行为不端或生活可责之人、以及不持守和教导正确信仰之人，已被我们拥有宗徒权威的先辈们剥夺了作为原告或证人的资格；我们也解除他们的这一职能，并在将来排除他们，免得我们本应坚守和拯救的人故意失足；并且，不仅（愿天主禁止！）所预言的审判将临到双方，而且我们也因他们的过失而丧亡（愿天主禁止！）。因为经上记载：\BibleQuote{他们立你为\Emph{宴席的主人}吗？你当照顾他们，好因他们而欢乐，并接受尊荣的冠冕，获得你所蒙拣选的认可。}\BibleRef{德 32:1}因为恶语影响人心，从心中发出善与恶、生与死这四样；而舌头不断的作用决定了它们。因此，上述人等应完全回避；在上述事项得到调查，并且查明他们清白之前，不应接纳他们：因为正确的祭献是遵守诫命，远离一切不义。\BibleQuote{远离邪恶是上主所喜悦的事，离弃不义是\Emph{赞颂的祭献}。}\BibleRef{德 35:1}因为经上记载：\BibleQuote{你要爱护你的朋友，对他忠诚。但你若泄露他的秘密，就不要再追随他。因为如同毁坏朋友的人一样，那丧失邻人友谊的人也是如此。如同从手中放走一只鸟，你放走了你的邻人，就不能再得回他。不要再追随他，因为他已远去；他像从网罗中逃走的羚羊，因为他的灵魂受了伤。你不能再束缚他，对受辱骂的人尚有和解，但出卖朋友的秘密，是可怜灵魂的绝望。以眼示意的人图谋恶事，没人能使他离开。当你面前时，他藐视自己的口，对你的言论表示惊奇；但最终他扭曲口舌，诽谤你的话语。我憎恨许多事，但不如憎恨他；上主也必憎恨他。谁向高处投石，石头必落在他自己头上；诡诈的打击必造成伤痕。谁挖陷阱，必掉进其中；谁在邻人的路上放置石头，必绊倒其上；谁为别人设网罗，必在其中丧亡。谁作恶，恶必临到他身上，他不知从何而来。傲慢产生嘲弄和辱骂；报复如狮子，等待他们。幸灾乐祸于义人跌倒的，必陷于网罗，痛苦在死亡之前吞灭他们。愤怒和忿怒都是可憎之物，罪人两者兼备。}\BibleRef{德 27:17}凡寻求报复的，必从主那里得到报复，而且他必保留自己的罪过。宽恕你邻人加给你的伤害，那么当你祈祷时，你的罪过也必得宽恕。一人怀恨另一人，而他竟向天主求伸冤？他对像自己一样的人不施仁慈，而他竟向至高者求赦免自己的罪过？他虽是血肉之躯，却怀有仇恨；他竟向天主求宽恕？谁将为他求情呢？记住你的终局，让仇恨止息；因为诫命面临朽坏和死亡。记住敬畏天主，不要对你的邻人怀有恶意。记住至高者的盟约，\OriginalTerm{轻视}{despice}你邻人的无知。远离纷争，你必减少你的罪过。\par

因为发怒的人挑起争端；有罪的人扰乱朋友，在和平的人中制造仇恨。因为树林中的树木如何，火焰也如何；人的力量如何，他的怒气也如何；他的财富如何，他使怒气上升的程度也如何。急促的争端会点燃火，急促的争吵会流血，作证（\Emph{testificans}）的舌头带来死亡。你若吹火星，它必如火烧起；你若唾其上，它必熄灭：二者都出自口中。愿谗言者和双舌之人受诅咒，因为这样的人使许多和平的人不安。第三条（\Emph{tertia}）舌头扰乱了许多人，将他们从一国驱赶到另一国；拆毁了富人的坚城，推翻了伟人的房屋。它颠覆了百姓的美德，毁灭了强盛的民族。第三条舌头赶走了诚实的妇女，夺去了她们劳动的果实。听从它的人永不得安宁，也永不能平静居住。鞭痕留在肉上，而舌头的打击折断骨头。许多人倒在刀剑之下，但不如死于舌头的人多。那能躲避恶舌的人有福了，他没有陷入它的怒气，没有负它的轭，也没有被它的绳索捆绑；因为它的轭是铁轭，它的绳索是铜索。它的死亡是最卑贱的死亡，坟墓比它更好。它的持续不能长久；但它必抓住不义者的道路，它的火焰不会燃烧义人。离弃主的人必陷于其中，它必在他们里面燃烧，永不熄灭；它必像狮子一样被派到他们身上，像豹子一样伤害他们。要用荆棘围住你的耳朵，不要听恶舌；为你的口安门，为你的耳加闩。熔炼（\Emph{confla}）你的金银，为你的话造天平，为你的口造正当的嚼环。当心免得你因舌头失足，跌倒在埋伏的敌人面前，你的跌倒无法挽回，直到死亡。\EditorialNote{德训篇第二十八章}\Quote{不要迟延归向主，不要一天推一天；因为祂的怒气忽然来到，在报应的时候祂必毁灭你。不要心向不义之财，因为在遮盖（即行刑，\Emph{obductionis}）和报应的日子，它们于你无益。不要随每一阵风摇摆，不要走每一条路；因为双舌的罪人就是这样被考验。要坚定地走在主的道上，在你的悟性的真理和知识中；愿和平与正义的道伴随你。听道要谦逊，好能明白；并用智慧给出正确的回答。你若明白，就回答邻人；若不明白，就用手捂口，免得你因愚昧的话被捉住，羞愧。智者的言语中有尊贵和荣誉；愚人的舌头是他的跌倒。不要被称为谗言者，不要被你的舌头捉住而羞愧。因为偷窃者将遭困惑和补赎，双舌者将受最坏的定罪。此外，谗言者招致憎恨、敌意和羞辱。要按公道对待小的和大的。}\BibleRef{德 5:7} 不要因邻人成为敌人。因为恶人必承受谴责和羞辱，每一个嫉妒和双舌的罪人也一样。不要像公牛一样在你心里自高自大，免得你的美德被愚蠢粉碎，它吞噬你的叶子，毁坏你的果实，你被留在旷野如枯树。因为邪恶的灵魂毁灭拥有它的人，使他被敌人嘲笑，并把他带入不虔者的命运。\BibleRef{德 6:1} 最亲爱的，要努力扶助被压迫的人，常常帮助穷乏的人；因为若有人救助受苦的弟兄、释放被俘者或安慰哀悼者，他应该毫不怀疑，这一切将被他所赐予的那一位报答，祂说：\Quote{你们对我这些最小兄弟中的一个所做的，就是对我做的。}\BibleRef{玛 25:40} 因此，要不断努力行善，使你们既在此世获得善工的果实，将来也享受天主的恩宠，以便日后配得进入天国。—— 在五月朔日前第四日（4月28日），最尊贵的塞维鲁和昆提安执政之年。\par

% structure: p0067 source=h2 target=section reason=relative-heading-rank; base=h2

\section{教宗\Person{安特鲁斯}书信}

\Emph{论主教与主教座的迁移}\par

致在贝提卡（\Place{贝提卡}）和托莱塔纳（\Place{托莱塔纳}）诸省被立为主教的最亲爱的弟兄们，主教\Person{安特鲁斯}在主内问安。\par

我最亲爱的弟兄们，我总愿收到你们真诚之爱与和平的喜讯，以便通过我们书信在你们中间的传播，彼此促进你们安康的迹象——如果那远古的仇敌给我们宁静，使我们脱离他的攻击；他从起初就是撒谎者（\BibleRef{若 8:44}），真理的仇敌，人的对手——为了欺骗人他先欺骗了自己——节操的对头，奢侈的主子。他以残酷为食，因节制而受罚；他憎恨斋戒，他的爪牙宣扬此类事为多余，毫无对未来的希望，并宣扬宗徒的那句话，其中说：\Quote{我们吃喝吧，因为明天就要死了。}（\BibleRef{格前 15:32}）啊，可悲的胆大妄为！啊，绝望心灵的诡诈！因为他唆使仇恨，驱散和睦。既然人心容易滑向更坏的部分，宁愿走宽路也不愿劳苦地走窄路，为此，最亲爱的弟兄们，你们要追随更好的，永远丢弃更坏的。行善，避恶，为要使你们被发现真是主的门徒。\par

如今，关于主教的调任，你们曾就此事希望咨询宗徒圣座，须知：这可合法地为公共福祉而行，或出于绝对必要，但不可随从任何个人的意愿或命令。我们神圣的导师、宗徒之长伯多禄，曾为公共福祉从安提约基雅调至罗马，以便他能在那里做更多的服务。欧瑟伯也因宗徒权威从某小城调至亚历山大里亚。类似地，斐理斯因其持守的教义和善生，经由主教、其他司铎及民众的共同同意，从他曾由市民选举被祝圣的城市调至厄弗所。因为那并非出于自己的倾向或野心之驱使，而是为公义或某种必要，经主要各方的商议并同意而被调任的人，不可被控以从一城游移至另一城之罪。也不可说他自小城调往大城，如果他被安置于此并非出于自利或自己的选择，而是或因受暴力逐出自己本座，或因某种必要所迫，且他不自傲而谦卑地由他人为地方或人民的益处而被调任并安置在那里。因为人看外貌，上主却看人心。且主借先知发言说：\Quote{上主知道人的思念，原是虚幻。}因此，那没有改变心意的人，并不改变其座位；那并非出于自己意愿，而是因他人的决定和选择而被改变的人，并不改变其城市。相应地，那并非为自己牟利或出于己意离开本城，而是如前所述，或被逐出本座，或受迫于某种必要，或基于司铎及人民的选举和命令而被调至他城的人，并不从一城游移至另一城。因为正如主教们有权按常规祝圣主教和其他等级的司铎，同样，每当任何有益或必要之事迫使他们时，他们也有权以上述方式调任和安置。既然你们在此事上询问我们的意见——虽然这些对你们并非未知——我们托付你们谨守这些，免得因某些人的无知，那更好、更有益的反被回避，而更无益的反被采纳，正如我们在神圣福音中所读：\BibleQuote{祸哉，你们经师和法利塞假善人！因为你们捐献十分之一的薄荷、茴香和莳萝，却忽略了法律上最重要的事：公义、怜悯与信德；这些是你们该做的，而其他也不可放弃。你们瞎眼的向导！滤出蚊蚋，却吞下骆驼。}\BibleRef{玛 23:23}合法的，\Emph{在他们那里}却是不合法；不合法的，反是合法。正如雅乃斯和杨布雷\BibleRef{弟后 3:8}曾抗拒真理，这些人心思败坏，爱享乐胜过爱天主，也教导：合法的便是非法，即主教应按前述方式从一城游移至另一城；而非法的事，他们却教导为合法，即向遭受困苦者施以怜悯是应省略的。也就是说，他们否认：应出于益处或必要，将属于他城的主教赐予没有主教并急需神圣主教牧职的人；以及应将另一个主教座席分配给遭受迫害或困苦的主教。他们也反对神圣圣经，因为圣经证述天主喜爱怜悯胜过审判。\par

请问，有什么比将一个人从无知的黑暗和缺乏经验的幽暗中拯救出来，并最终通过真信仰教义的滋养，不是为获利或野心，而是为教导和建立而恢复他，更大的爱德或更有效的虔敬服务呢？[因为他成为了，可以说是，残者的手，跛者的脚，盲者的眼\BibleRef{约 29:15}，那为被无知的黑暗所包裹的人打开智慧和知识的宝库，并为这样的人敞开光明之灿烂和上主道路的人。]\par

如今，对于那些忍受天主言语之饥荒的人，以及那些为公共利益而被安置在其他城市、忍受困苦的主教们，所显示的仁慈并非微小。而那些否认这点的人，虽然具有敬虔的外貌，却否认了它的能力。\\BibleRef{弟后3:5} 在这样的事上，我不承认种族（\\OriginalTerm{种族}{prosapiam}）。然而，如果哪位明智者，因这场风暴（或时节）的紧迫而与愚昧者中的其他领袖联合，并因参与他们的行为而被玷污，但明智者的卓越，即使他可能偶然知晓他们的过错，也使他不能将自己作为罪人的领袖。公共利益与需要的动机是一回事，而自私、僭越或私人偏好的动机是另一回事。因自私、僭越或私人偏好，主教不得从一个城市调往另一个城市，只能因公共利益和需要才可以。这是无人否认的事，除非那些论他们所说的：\\BibleQuote{他们因酒而失迷；他们不识先见者；他们不知审判。} \\BibleRef{依28:7} 因为如果我被迫叙述那些已经终结的事，我将向你们显示，比较这些行为并不能带来任何安慰。但是，最亲爱的诸位，\\BibleQuote{你们应站在路上，观看，询问主的古道，看哪是善道正路，行在其中，你们必得灵魂的安息。} \\BibleRef{耶6:16} 并且，按照智慧之言来说：\\BibleQuote{爱慕正义，你们审判大地的人！应以善意思念主，以纯朴的心寻求祂。因为，祂不被试探的人所寻获，却向不怀疑祂的人显示自己。因为邪僻的思想使人远离天主；愚人一受试探，便显露自己的无知。因为智慧不进入恶意的灵魂，也不住在一个屈服于罪恶的身体内。因为圣神必远避欺诈，远离无知的思念，不义一到，即不复存在。因为智慧是慈善的神，决不赦免渎神者的口舌。因为天主是他内心的见证，是人心肺腑的真探察者，是他唇舌的听闻者。因为主的神充满了世界，包罗万象，洞悉一切语言。为此，出言不义的人，绝不能隐藏；惩罚的裁判，也绝不放过他。因为不敬神者的计谋，必受调查；他的话，必然传到主面前，作为他罪恶的凭证。嫉妒的耳朵，能听万事；怨声四起，也隐藏不住。为此，你们要谨防无益的怨言，禁止口舌的毁谤；因为暗中的言语，不会空虚无用；说谎的口舌，杀害灵魂。不要因生活腐化而自招死亡，也不要因双手作恶而自取灭亡。因为天主并未创造死亡，也不喜欢生灵灭亡。祂创造了万物，为叫它们生存；世上的生物，都有生命力，本身没有致命的毒素；阴府在地上也没有权势。正义是永存不朽的，不义却招致死亡。不敬神的人，用言语用手势召唤了死亡；他们以为死亡是朋友，便沉溺于其中，并与它立约；因为他们配作死亡的伙伴。} \\BibleQuote{因为他们心中思念，但不正确：我们的生命短促而烦恼；人死时，没有补救之法，也没有人从阴府回来。因为我们是由虚无而生，以后也将好像未曾存在过。我们鼻中的气息是烟雾，言语是一点火星，心房跳动即告熄灭，身体化为灰烬，灵魂消散如清风。我们的生命如云影消逝，如雾气被太阳驱散、被热量压倒。我们的名字随时间而被遗忘，没有人记得我们的工作。我们的时间如影子消逝，死后永无复返；因为印已封好，无人再能回来。} \\BibleRef{智2:1} 为此，每个人必须谨慎保守自己，为自身的益处而自我警醒，好使当他的末日和生命的终结来临时，他不至堕入永死，而获得永生。因为受我们管辖之人的行为由我们审判，但我们自己的行为却由天主审判。此外，有时主教因民众的过失而变坏，以致跟随他们的人跌得更深。当首脑衰弱时，身体的其他肢体也受影响。败坏善人生命与道德的人，比掠夺他人财产和货物的人更为卑劣。每人要注意，既不要舌头痒，也不要耳朵痒；也就是说，既不要自己毁谤他人，也不要听别人毁谤。他说：\\BibleQuote{你坐着，议论你的兄弟；诽谤你母亲的儿子。} 每个人要远离毁谤的舌头，谨守自己的言语，并明白他们谈论他人的一切，都将成为审判他们自己的依据。没有人乐意对不愿听的听众说话。最亲爱的诸位，你们要关心不仅保持眼睛纯洁，也要保持舌头纯洁。不要因你们而让另一个家庭知道某人家中发生的事。愿所有人都拥有鸽子的纯朴，不对任何人设诡计；也拥有蛇的机警，以免被他人的诡计所欺骗。评判他人或在教会服务人员身上说任何不利的话，不属于我卑微的身份和尺度。我绝不说任何不利于那些继承宗徒地位、用他们神圣的口舌祝圣基督身体的人的话；通过他们，我们也成了基督徒，他们拥有天国的钥匙，并在审判之日以前施行审判。此外，古代法律中记载：凡不服从司祭的，要么被民众在营外以石头砸死，要么伏首刀下，以血抵偿他的僭越行为。\\EditorialNote{申命纪第十七章}\par

如今，叛逆者却被属神的惩罚所剪除；被逐出教会后，被恶鬼的狂口撕裂。因为以天主为产业的人，理应脱离世上一切妨碍，事奉天主，好使他们能说：\Quote{主是我的产业的分。}\BibleQuote{主啊！祢的圣神在一切事上何其良善、可亲！}\BibleRef{智 12:1}祢宽恕一切，因为一切都是祢的，主啊，祢爱护灵魂。\Emph{因此，祢渐渐责罚那些犯罪的人}，警告他们得罪之处，劝诫他们，使他们离开邪恶，相信祢，主啊。\BibleRef{智 12:2}\BibleQuote{但是祢，我们的天主，是宽仁的、真实的、坚忍的，以慈爱治理万有。因为我们若犯罪，仍属祢，知道祢的权威；若不犯罪，也知道我们算是祢的。}\BibleRef{智 15:1}\BibleQuote{敬畏主的人的精神，主必眷顾；在天主眼中，他们是有福的。}\BibleRef{德 34:13}为此，最亲爱的弟兄们，\BibleQuote{一切坏话都不可出你们的口，但只说有益于建立人的话，使听者获得恩宠。不要叫天主的圣神忧郁，因为你们是在祂内受了印记，以待得救的日子。一切毒辣、怨恨、忿怒、争吵、毁谤，连同一切恶毒，都要从你们中除掉。你们要彼此以仁慈相待，要心慈，互相宽恕，如同天主在基督内宽恕了你们一样。}\BibleRef{弗 4:29}\BibleQuote{因此，你们要效法天主，如同蒙宠爱的儿女一样；并要在爱德中行事，就如基督爱了我们，为我们舍弃了自己，当作馨香的祭品和牺牲，献给天主。至于邪淫、一切污秽或贪婪，在你们中连提也不可，这样才合乎圣徒的身份；更不可说淫秽、愚蠢或轻薄的话，这都是不相宜的；反而要感恩。因为你们确知：凡是淫荡的、不洁的、贪婪的（即偶像崇拜者），在基督和天主的国里，都得不到产业。不要让人以虚言欺骗你们，因为正是为了这些事，天主的义怒才降在叛逆之子身上。所以，你们不要同他们分享。从前你们原是黑暗，但现在你们在主内却是光明；你们要像光明之子一样行事（光明所结的果实就是良善、正义和诚实），时时寻求上主所喜悦的事。不要参与黑暗无益的行为，反而要指摘它们。因为他们暗中所行的事，连提说也是可耻的。一切事被指摘，便由光明显露出来，因为凡显露出来的（\Emph{manifestatur}）就是光明。因此经上说：\InnerQuote{你这睡眠的，醒来，从死者中起来，基督必要光照你。}所以，弟兄们，你们要小心行事，不要像愚昧人，却要像有智慧的人，善用时机，因为时日险恶。为此，不要作糊涂人，但要明白什么是主的旨意。不要醉酒，醉酒使人丧德；却要充满圣神，用圣咏、诗词及属神的歌曲，彼此对谈，在心中歌颂上主；时时为一切事，因我们的主耶稣基督之名，感谢天主父；要怀着敬畏基督的心，彼此服从。}\BibleRef{弗 5:1}所以，弟兄们，你们要站稳，持守宗徒和宗座的传统，\BibleQuote{愿我们的主耶稣基督和那爱了我们、并藉恩宠赐给我们永远安慰和美好希望的天父，安慰你们的心，并使你们在一切善行善言上得以坚固。}\BibleRef{得后 2:15}\BibleQuote{最后，弟兄们，请为我们祈祷，愿主的圣言快快传开，并受到光荣，如同在你们中一样；也愿我们能摆脱无理和邪恶的人，因为不是人人都有信德。但主是信实的，祂必坚固你们，保护你们免受邪恶。}\BibleRef{得后 3:1}因此，要常常把心放在天主的德能（\Emph{virtute}）上，常常抵抗邪恶，并按照先知的话，\Quote{传告后代；因为这位天主是我们永远的天主，祂要统治我们直到永远。}因此，你们这些为主立在守望台上（\Emph{in specula}）的人，务必阻止和拦阻那些对弟兄图谋诡计、或激起叛乱和诽谤的人。因为用言语欺骗人是容易的，但在主面前却不然。所以你们应当斥责这种人，远离他们，为要使这一切黑暗彻底消除，晨星照耀他们，喜乐在他们心中升起。\BibleQuote{\Emph{我们}在主内深信你们，弟兄们，你们现在遵行、将来也必遵行我们所吩咐的事。}\BibleRef{得后 3:4}因为你们越向他们显出仁慈，就越能从他们所事奉的全能天主那里得到更大的回报。愿全能的天主保护你们，赐你们持守尊荣和诫命；愿光荣和尊威归于全能的天主圣父，及祂的独生子我们的救主，同圣神，直到永远。阿们。\par

颁于四月朔日前第十二日（即三月二十一日），最尊贵的马克西米安和阿弗里卡努斯执政之年。\par

% structure: p0076 source=h2 target=section reason=relative-heading-rank; base=h2

\section{法比安教宗第一封书信}

\Emph{致公教会全体神职人员。论不应允许为自己辩护者，以及不与受绝罚者交往的责任。}\par

致各地公教会中亲爱的主内弟兄们，法比安在主内问安。\par

我们受天主诫命和宗徒规诫的劝谕，要以不懈的关心为所有教会的位置而警醒。因此，你们应当知道罗马教会神圣事务中所行之事，以便效法她的榜样，成为那被称为你们母亲之教会的真正子女。因此，我们按照从祖先所接受的制度，在罗马城内维持七位执事，分驻于该邦的七个区域，他们每周轮值所任的职务，在主的日子和庄严庆节，与副执事、辅祭者及以下各等级的服务者一同执行，并随时准备履行宗教义务和所有托付给他们的事。同样，你们也应当在各自不同的城市中，依便利而行，以便热忱而有规律地履行宗教义务，毫无迟延或疏忽。此外，我们同样规定了七位副执事，他们应\OriginalTerm{侍立}{imminerent}于七位公证人旁边，将殉道者的历史汇集成一份完整而准确的记录，提交给我们审阅。我们也督促你们众人这样做，免得在后世对这些事产生疑惑或争议。\BibleQuote{因为凡所写的，都是为教训我们而写的。}\BibleRef{罗 15:4}凡我们时代真实书写的事，都是为教训后世。因此，我们嘱托这些职责须交托给最忠信的人，使其中毫无虚假，以免（愿天主禁止）从其中生出对信徒的冒犯。为此，我们也以慈父的善意恳求你们的爱德，愿圣教会如今在一切事上得见你们爱德的美意，并在有需要时获得你们恩惠的安慰。既然你们热忱的良善使我们确信，不应在何事上不信任你们，反而应当以完全的信赖将这些事托付给你们，如同托付给我们教会的智慧之子；那么，在时机不重要的情况下，你们的德行更应当奋发努力，竭尽一切可能，以全部热忱驱除责难。我们也按照宗徒的话劝勉你们，要\BibleQuote{坚定不移，常常在主的工作上充盈，因为知道你们的劳苦在主内不是徒然的。}\BibleRef{格前 15:58}又在另一处说：\BibleQuote{你们要警醒，祈祷，在信德上站稳。要作大丈夫，要刚强。一切事都当以爱德而行。}\BibleRef{格前 16:13}此外，我们愿你们知道，在我们时代，因我们的罪孽困扰，那古时的仇敌，常像咆哮的狮子巡游，寻找可吞食的人，\BibleRef{伯前 5:8}挑起了他，\Person{诺瓦图斯}从\Place{非洲}上来，将\Person{诺瓦提安}和基督的某些其他宣认者从基督的教会中分离，说服他们接受恶毒的教义。弟兄们，要远离这样的人，并提防所有持守并教导与宗徒及其继承人所持守并教导不同的信仰和教义的人，免得（愿天主禁止）跟随他而落入撒殚的罗网，被他的锁链束缚。因此，我们以最恳切的祈祷恳求你们的兄弟之爱，愿你们在圣洁的祈祷中记念我们卑微之人，恳求和祈求上主的天主，使我们以及基督的圣教会——我们的母亲，以祂宝血赎回的——能从埋伏我们的撒殚罗网中，并从烦扰和邪恶的人中被拯救出来，使天主的圣言得以自由传扬并受光荣，并使他们以及所有教导悖逆真理之人的恶毒教义被推翻而消灭。我们也恳求你们在虔诚的祈祷中热切祈求，愿我们的天主和主耶稣基督——祂愿意所有的人都得救，无人丧亡，\BibleRef{弟前 2:4}——以祂无限的全能，使他们的心回转于纯正的教义和公教信仰，以便他们能从魔鬼的罗网中被解救，那些被牠掳掠的人，而与我们的母亲教会的子女合而为一。也要记念你们的弟兄，怜悯他们，尽你们所能的一切方式为他们劳苦，使他们不致丧亡，而是因你们的祈祷和你们良善的其他努力，在主内得救。因此，在这些事上你们要如此行事，好使你们证明自己是天主圣教会顺从而忠信的子女，并获得赏报。这些人以及所有不教导真道、不持守真信仰的人，不能作任何真信徒的控告者，因为他们被烙上了恶名，并被宗徒之剑从我们圣母亲教会的怀抱中切断，直到他们归正于正确的行为和信仰。因此，凭藉宗座的权威，并与同一宗徒和普世教会的所有儿子一致，我们决定：凡在公教信仰上受嫌疑的人，都不能被接纳为持守真信经之人的控告者；因为嫌疑总应当排除。故此，凡由在真信仰上受嫌疑之人所提出的控告，理当被驳回。凡不认识天主圣三信仰之人，也完全不可信赖。同样，我们摒弃并排除所有那些被历代圣父教令所咒骂之人，使其不得参与控告信徒。因此，信者应常与不信者分别，义人与不义者分别；因为不信者和恶念之人，总是尽一切可能搅扰信者，力图毁灭他们；故此他们不应被接纳，而应被拒绝，并完全远离，以免他们毁灭或诽谤信者。为此，亲爱的诸位，要提防这等人的陷阱，我们知道有许多人已落入其中。要提防这等人的圈套（或言飞镖），以及古仇敌的计谋，我们曾看见连与我们亲密的人都在我们面前受伤跌倒。要警戒埋伏者的活结，他们惯用以勒死同伴和同志。不要跟随这样的人，而要远离他们。要按照真理之声，\BibleQuote{机警如同蛇，纯朴如同鸽子。}\BibleRef{玛 10:16}注意，不要徒然奔跑或劳苦；而要彼此以祈祷和恳求扶持，努力遵行天主的旨意；并从我提及的那些人，若他们显示自己不可悔改，就在一切事上远离他们。同样，也要远离所有那些宗徒提及的人，他说：\BibleQuote{与这样的人，连饭也不可同吃。}\BibleRef{格前 5:11}因为这些人，如同前者，当被拒绝，在未向教会作出满意前不得被接纳。因为那些不许与之吃饭的人，显然是与其余弟兄的一切交往隔绝，直到作出满意。因此，他们不应也不能被接纳为控告信徒之事的提出者，而应被排除于其团体之外，直到上述满意作出，免得这些人也变得像他们一样，或受他们的绝罚所牵连；因为宗徒们为此已有谕令，说：与绝罚者不可有共融。若有人明知故犯地忽视规则，与绝罚者在家中唱诗，或与他们交谈或祈祷，那人将被剥夺共融的特权。故此，这样的人在一切事上都当提防，不得接纳，因为按照宗徒，不单是行这类事的人被定罪，连那些同意他们的人也是如此。\BibleRef{罗 1:32}因此，蒙福的宗徒之长伯多禄，在祝圣\Person{克莱孟}时向百姓说话，也说了这样的话：若这位\Person{克莱孟}因行为而敌视某人，你们不要等他直接对你们说：不要与此人友好。但你们应当留意他的心意，如合宜一样，并无需直接吩咐就附和他；并要远离你们察觉他所敌视的那个人，与他所不交谈的人不要交谈，以便每个有过失的人，既然渴望得到你们所有人的友谊，就会热切地尽快与那管辖众人者修好，从而借此\OriginalTerm{恢复灵性的健康}{redeat et salutem}，当他开始服从首长的命令时。\par

然而，若有人不友善，与他（他的首领）不与交谈的人交谈，这样的人属于那些试图消灭天主教会的人；虽然他外表似乎与你们同在，但内心和意念却反对你们。这样的人比那些外在、公开敌对的人是更危险的敌人。因为此人以友谊为伪装，扮演敌人的角色，分散并毁坏教会。因此，亲爱的诸位，在这些宗徒的教导中，我们警告并教导你们，你们的爱德在其中受到教导（\OriginalTerm{更加确实}{effecta certior}），日后可能更加谨慎和明智地行事，使悖逆和不忠之人无力伤害忠实和良善之人；因为这样的人和所有不虔敬之人的希望，如同被风吹散的尘土；如同被风暴驱散的薄沫；如同被暴风四散的烟雾，又如同只停留一天的旅客的记忆。\BibleRef{智 5:14} 亲爱的诸位，要极其小心地提防、躲避并拒绝这样的人，如果他们显出伤害之意。因为世俗的法律与教会的法律一样，都不接纳伤害者，而是拒绝他们。因此经上记载：\Quote{恶人的口吞吃不义。} \BibleRef{箴 19:28} 主借先知说：\Quote{与圣洁者，你将显为圣洁；与乖戾者，你将显为乖戾；与卓越者，你将显示自己卓越（\OriginalTerm{被拣选者}{electus}）；与无辜者，你将显为无辜。} 宗徒说：\Quote{不良的交往败坏善行。} \BibleRef{格前 15:33} 因此，如已指出的，恶人总要被躲避和远离，而良善和正直的人要坚定地追随，以便我们尽可能避免懈怠的危险。为使这瘟疫不至更广蔓延，我们要以最严厉的态度将其从我们中间割除；因为哪里有虔敬的勤勉，哪里就不会有鲁莽的冒失。你们每个人，靠这宗徒的代表权，按自己的力量行事，在兄弟之爱和虔诚的敬虔中，努力端正自己的品行，互相帮助，常守爱德，毫不间断地持守天主的旨意，使我们一同赞美主，永远不知疲倦地感谢祂。亲爱的诸位，愿你们在主内安好，并靠主的帮助尽力完成上述诸事。\Signature{发于七月一日，最显赫的马克西米努斯（或马克西穆斯）与阿非利加努斯执政之年。}\par

% structure: p0081 source=h2 target=section reason=relative-heading-rank; base=h2

\section{法比盎教宗第二封书信}

\Emph{致所有东方主教：圣油应每年祝圣更新，旧油应在教堂焚烧；亦论控告司铎之事，以及羊群不得妄责牧者，除非牧者在信仰上犯错。}\par

法比盎，罗马城主教，致所有东方主教及全体信友，在主内问安。\par

你们对宗座之位的热爱，要求我们既不能也不应拒绝给予劝勉。此外，我们的前任为许多地区的主教做了这事；而兄弟之爱和服从的义务也使我们这些因天主的丰厚恩慈而居于同一座位的人负有同样的责任。因此，你们的关怀应当留意，不可让懈怠导致疏忽，也不可让冒失扰乱那些由宗徒及其继任者所规定、并在圣神的默感下所确立的事。既然正确秩序的使用要求定义某事，那么已经定义的事就不应被违反。\par

% structure: p0085 source=h3 target=subsection reason=relative-heading-rank; base=h2

\subsection{一、每年应制作新圣油，旧油应焚烧。}

如今，在你们信函所提及的其他事项中，我们发现一些你们辖区的主教采取了与你们及我们不同的做法：他们不在每年的主的晚餐上制备圣油，而是将其保存使用两三年，一次制备足量的圣油。因为他们声称（如我们在该信函中所见），无法每年获得香脂；而且即便能够获得，也无需每年制备圣油，只要一次制备的圣油足够多，便不需再做新的。然而，如此想法的人实属错误；他们的言论更像是疯癫者而非清醒者。因为在那一天，主耶稣在与门徒共进晚餐、并为他们洗脚之后，按照我们的前辈从圣宗徒那里领受并传留给我们的传统，教导他们制备圣油。洗脚象征我们的圣洗，而圣洗藉圣油的傅油得以完成和坚振。正如那一日的隆重庆典应每年举行，圣油的制备也应每年进行，且需逐年更新并分施给信友。因为\Emph{这新圣事的原料}应每年更新，并在上述之日进行；旧的圣油应在圣堂中焚烧。这些事我们从圣宗徒及其继承人那里领受，并托付给你们保管。罗马圣教会和安提约基雅教会自宗徒时代起便守护这些事；耶路撒冷和厄弗所的教会也持守它们。在这些教会中主理的宗徒教导了这些事，并规定旧的圣油应予焚烧，允许其使用不超过一年，之后命令他们使用新的而非旧的原料。因此，若有人胆敢违抗这些事，他应明白，宽恕之门已由你们及所有正直之人向他关闭：因极堕落心性的邪僻教义，在过于放纵地使用缰绳时，滑入僭越之罪；除非得到明智者的全力支持和纠正，它绝无可能被驱逐。至于圣教会普世一致遵守的、关于神圣奥迹和洗礼对象的所有惯例，不应漠不关心，以免为无谓的努力和迷信开道。因此，我们不应将信友中未经教导的心灵引入我们前述的做法，因为他们应受教导而非被操纵。因为善行带给我们幸福，恶行则以悲伤之刺刺痛我们。但在此处，无论我们身处何地，我们都在强盗之手和狂怒之狼的獠牙之间，悖逆之人占据了真羊的位置。唯有犬吠和牧人的牧杖才能遏制狼的狂怒。此外，那些无法用药治愈的伤口，必须用刀切除。我们也不能保持沉默，因为在此处试图将一些人从非法之事唤回时，我们受本分本能的驱使，被主安置在瞭望台上，目的就是要通过制止应禁止之事、决定应遵守之事，来证明我们警觉的勤勉。\par

% structure: p0087 source=h3 target=subsection reason=relative-heading-rank; base=h2

\subsection{二、论主教不受诽谤之控告或伤害的权利}

你们也愿意就控告司铎一事向我们请教，正如我们在你们上述信函中所见——这件事，我们从同一封信中也得知在你们当中极其频繁。此外，你们还提到，许多人注意到，不少身居教会尊位者的生活与其所用以服务百姓的言论和圣事不相符。可怜的人哪！他们注视这些时，竟忘记了基督；基督很久以前就告诉我们，应当遵守天主的律法，而不应效法那些言行不一的人；祂甚至容忍了叛徒直到最后，还派他与其余的人一同宣讲福音。因为宗徒们没有这样的习俗，也没有教导这是可取的习俗。同样，他们的继承人，藉着天主的神预见将来，也在这类事上作了大量决定。此外，正如你们在《宗徒大事录》中读到：\Quote{那时，众信徒都一心一意，没有人说自己的财物是自己的，而是凡物公用。}\BibleRef{宗 4:32}因为他们当中没有彼此控告的事，除了友善的；在跟随他们的人或信徒当中，也永远不应有这类事。因为主说：\Quote{己所不欲，勿施于人。}祂又说：\Quote{你应当爱近人如你自己；}以及：\Quote{爱不加害于人。}\BibleRef{罗 13:10}据此，宗徒们和他们的继承人古时已规定：有嫌疑者，或昨日、前日、不久之前曾是仇敌者（因他们招致嫌疑），或品行不善者，或其生活可指责者，或其信仰真道可疑者，不得被接纳提出控告。同样，对于信仰、生活和自由不明者，或带有恶名污点者，或陷入罪过罗网者，也如此裁定。此外，那些自身不能合法成为司铎、且不属于他们品级的人，既无权也无能力控告司铎或圣职人员；因为正如司铎和其他圣职品级的人被禁止控告平信徒，平信徒也应当被禁止和排除在控告前者的权利之外。正如前者不应被后者钦佩，后者也不应被前者钦佩：因为主的司铎的生活本应与这些人的生活有别，在诉讼的事上也应与他们分离；\Quote{因为主的仆人不应当争辩。}\BibleRef{弟前 2:24}最亲爱的弟兄们，你们要尽最大能力禁止这样的控告，以及一切不义和有害的竞争，因为争辩必须全然避免。\Quote{因为义人一日七次跌倒，仍能起来；但恶人必陷于祸患。你的仇敌跌倒时，不要欢喜，}撒罗满说，\Quote{他失足时，你心里不要快乐；免得主看见而不悦，就转开祂的怒气离开他。不要因作恶的人而烦恼，也不要嫉妒恶人；因为恶人没有将来的盼望，恶人的灯必要熄灭。不要嫉妒恶人，也不愿与他们同在；因为他们的心图谋强暴，他们的嘴唇说欺诈的话。}最亲爱的，要提防这些事。要默想这些事，在一切事上安慰弟兄们；因为正如真理亲自所说：\Quote{如果你们彼此相爱，众人因此就可认出你们是我的门徒。}\BibleRef{若 13:35}因为如果在世俗事务中，各人的权利和适当位置都为其保留，那么教会秩序的事岂不更应避免混乱！这个权利若要得到恰当遵守，就不可一味看重权力，而要一切遵循公平。因此，这是一项既定的责任：各个区域的教区主教应当小心照顾所有在其管辖下生活的人，并怀着对天主的敬畏处置所有涉及他们的事和与他们相关的一切事务。因此，若有\Emph{主教}忽略自己管区的事，而插手他人的事，是极不公平的。但那些负责祝圣这样的人为司铎、且已祝圣他们的人，应当以适当而正常的行政管理来安排这些人的生活与判断；因为正如律法所说：\Quote{凡移动邻舍地界的人，是可咒骂的。众百姓都说：阿们。}\BibleRef{申 27:17}为此，弟兄们，天主早已预定了你们和所有握有最高司祭职的人，要你们除去一切不义，剪除傲慢，帮助那些在司铎职中劳苦的人，不给他们受辱和麻烦的机会，而要给遭受毁谤和侮辱的人提供援助，剪除施行毁谤和侮辱的人，并为了主在祂的司铎中而帮助主。此外，主为自己拣选了司铎，使他们向祂献祭，向他们的主奉献礼品。祂也命令肋未人受他们管束，服务于他们的职务。为此，祂向梅瑟这样说：\Quote{亚郎司祭的儿子厄肋阿匝尔应为肋未人的首领，监察那些看守圣所的人。}\BibleRef{户 3:32}因为关于这些人，主这样对梅瑟说：\Quote{你要以肋未人代替以色列人中所有头胎的，以肋未人的牲畜代替他们的牲畜；肋未人要归我：我是上主。}\BibleRef{户 3:45}如果主愿肋未人归于自己，那么祂岂不更将司铎归于自己！关于这些人，祂说：\Quote{凡外人走近，应处死。}\BibleRef{户 1:51}此外，凡属于主的事物都应当谨慎处理，不可轻易损害；因为即使在人当中，那些忠心料理主人事物、忠诚待之、正确遵守主人命令而不违犯的人，也被视为忠信。另一方面，那些粗心疏忽料理主人事物、藐视命令、不按当尽之责遵守的人，则被视为不忠信。\newline{}\par

因此，我们将这些事摆在你们面前，为使那些如今尚不知晓的人得以知晓：即那些从众人中被主选召、并蒙祂意愿归为己有的司铎，也不应被轻率对待，不应受伤害，也不应被草率控告或斥责，除非由他们的上司，因为主已选择将他们的案件留给自己，并按祂自己的审判施行报应。因为在这些及其他主的诫命中，信众得以区分，同时不信者受到责备。因为这些人与其被信众作为受\OriginalTerm{斥责}{exprobrandi}的对象，不如被他们容忍；就如麦子中有糠秕直到最后簸扬，也如好鱼中有坏鱼直到他们被分开——这分开将要在岸上，即世界末日。所以，那个人绝不能通过人的审判断罪，因为天主已将他保留为自己的审判，使天主的旨意——祂曾决意拯救那已丧亡的——不可更改。因此，既然祂的旨意不容改变，人不可僭越那些未赋予他的事。而这就是宗徒所说的那句话的含义：\BibleQuote{你们彼此告状，已经是你们的大错了。为什么不情愿受欺呢？为什么不情愿吃亏呢？}\BibleRef{格前 6:7}为此，我们主的话也可指涉：\BibleQuote{若有人要拿你的外衣，并要控告你，连你的内衣也任他拿去。}\BibleRef{玛 5:40}在另一处：\BibleQuote{那拿走你东西的，不要再向他索回。}\BibleRef{路 6:30}此外，有些事若未在圣经中显示其更严重的含义，可能被视为微不足道。若非真理亲自告诉我们，有谁会认为那对自己的兄弟说\BibleQuote{傻子}的人该受地狱之火呢？\BibleRef{玛 5:22}此外，那些犯下宗徒所说的\BibleQuote{行这样事的人，不能继承天主的国}\BibleRef{迦 5:21}的罪的人，必须严加防范；若他们不自愿寻求改正，就应该强制他们，因为他们被烙上耻辱的污点，并坠入深渊，除非藉司铎的权威给他们带来援助。那些他曾说\BibleQuote{这样的人，连和他一起吃饭也不可}\BibleRef{格前 5:11}的人，也要同样对待；因为这样的人被烙上耻辱，直到藉司铎的权威得以恢复，并重新纳入我们神圣慈母教会的怀抱；因为那些在我们之外的人不能与我们相通。显然，这些人在我们之外，应当与我们分离，因为与他们一同进食或共食是不合法的。同样，凡犯有任何形式的卑鄙和不名誉之罪的人，都被视为可耻；凡武装起来对抗父辈的人，也被视为可耻。\BibleQuote{沙盐、铁堆，比一个无智、又愚妄又邪恶的人，更容易负荷。}\BibleRef{德 22:15}\BibleQuote{缺乏智慧的人思想虚妄的事；愚昧而迷途的人幻想愚妄的事。}\BibleRef{德 16:23}因为他们的猜疑曾使许多人倾覆，他们的意见使他们陷于空虚。\Quote{顽梗的心终必遭祸；喜爱危险的人必在其中丧亡。一颗心分走两条路，必不得安宁；邪恶的心在其中必跌倒。恶人的心必为忧愁所负累；罪人必罪上加罪。}圣宗徒们和他们的继承人，心中记念这些事，并因充满圣神而预见恶人的行径，且顾念纯朴的人，决定控告司铎应是一件困难的事，或者决不进行，以免他们被恶人摧毁或取代。因为若这件事对世俗和恶人变得容易，就没有人留下，或只有极少数；因为过去和现在一向如此——更糟的是，日益加剧——恶人迫害善人，属血肉的人敌视属神的人。因此，如前所述，他们规定这样的人根本不应被控告；若无法避免，对这样的人的控告应极为困难。他们还决定，如上所述，哪些人不应承担那职务；并进一步决定，主教不应被逐出他们自己的专属席位和教会。但若在任何情况下，在他们合法的席位及所有财产依那些法律得到归还之前，那\Emph{控告}之事被着手进行，他们绝不应被任何人控告或指控，也不应对此类指控回答任何人，除非他们自愿这样做。但在他们恢复原位之后，如前所述，且藉那些法律将全部财产归还给他们，当他们的各项事务安顿有序之后，他们应有较长时期处理案件；此后，若有必要，应正式传唤他们，让他们出庭；若事情看来合理，他们应在弟兄们的帮助下回答控告者的陈述。因为只要他们的财产，或他们的教会和权益，被敌人或任何人占有，就没有任何理由允许对他们提出任何指控。而且，只要他们被剥夺了他们的教会、财产或权力，无论上级或下级，都不得以任何方式自由地对他们提出任何指控。同样，也已规定，我们也确认同样的法令，并在此规定：如果任何圣职人员成为他的主教的敌人或诋毁者，并试图指控他们，或密谋反对他们，应在司法调查考虑之前，立即将他从圣职行列中撤除，并交给\OriginalTerm{法庭}{curiæ}，他应毕生热心服役于该法庭，且保持可耻，没有任何恢复的希望。任何人都不得同时担任控告者、法官或证人；因为在每一项司法调查中，总须有四人出席：即选任的法官、控告者、辩护人和证人。同样，我们以宗座权威规定并命令：羊群不应胆敢控告他们的牧人，即他们被托付给其照顾的人，除非他在信仰上犯错；因为上级的行为不应被口的剑击打；门徒也不能在师傅之上，正如真理的声音所说：\BibleQuote{门徒不能高过师傅，仆人也不能高过主人。}\BibleRef{玛 10:24}\par

在天主及人面前，骄傲可憎，一切不义可恶。\Quote{上主消灭了骄傲人的记念，却保留了谦卑人的记念。人的后裔将受尊敬，那敬畏天主的\Emph{后裔}必受尊敬。那违背上主诫命的后裔，必受羞辱。在弟兄中，为首的是可敬的；敬畏上主的人，在他眼中是尊贵的。我儿，撒罗满说，你要以温良保存你的灵魂，并要给予那应受尊敬的人以尊敬。} \Quote{未查明以前，不可责斥任何人；查明以后，才可公正地责备他。未听见原因以前，不可回答一句话；也不可在年长者中间插嘴谈论。} \BibleRef{德 11:7} 他们效法诺厄的儿子含的榜样，公开宣扬父亲的过失，或妄加控告和毁谤，正如含那样，他没有遮盖父亲诺厄的羞耻，反而将其暴露出来，以供嘲笑。同样，那些因循闪和耶斐特榜样的人，被证明为义理，因为他们虔诚地遮盖，不寻求揭露在父亲们身上发现的过错。因为如果一位主教的信仰偶然出错，他首先应由他的属下（\OriginalTerm{由他的属下}{a subditis suis}）私下加以纠正。如果他显明是不可救药的（但愿天主阻止），那么就应在省区主管或宗座面前控告他。至于他的其他行为，他的羊群和属下宁可容忍他，也不应控告他或公开毁谤他；因为当属下在这些事上犯错时，就是违抗了那将他们置于他面前的天主的命令，正如宗徒所说：\Quote{谁反抗权柄，就是反抗天主的命令。} \BibleRef{罗 13:2} 但敬畏全能天主的人，决不会做任何违背福音、违背宗徒、违背先知或圣教父的传统的事。因此，司祭应当受到尊敬，而不应受伤害或斥责。我们在德训篇中读到：\Quote{你当全心敬畏上主，尊敬他的司祭。你当全力爱慕造你的主，不可离弃他的臣仆。你当全心尊敬天主，并尊敬司祭，且预先用双肩洁净（\OriginalTerm{用双肩洁净}{propurga te cum brachiis}）自己。你当将初熟之物，按所吩咐的给他；并以少量之物洁净你的疏忽。你应将你双肩的礼物，圣化的祭品，以及圣物的初熟之物献给上主。你当向穷人伸手，使你的赎罪和祝福得以成全。} \BibleRef{德 7:29} 我们希望这些事不只为你们所知，而且通过你们为所有弟兄所知，使我们能同心合意地居住在基督内，不因争吵或虚荣而自高自大，不求人的喜悦，只求天主我们的救主的喜悦。尊荣与光荣归于他，直到永远。阿们。\par

% structure: p0091 source=h2 target=section reason=relative-heading-rank; base=h2

\section{教宗法比盎第三封书信}

\Emph{致希拉略主教：论应拒绝外来的裁判，被告应在本地进行诉讼；凡提出控告者，应书面表示其证明能力，若不能证明其所指控，应承担其所提出的刑罚。}\par

\Person{法比盎}，致我至爱的弟兄\Person{希拉略}主教。\par

我们应思念天主对我们的恩宠，祂出于祂慈爱的眷顾，为此原因将我们提拔至司祭职的高峰，使我们通过坚守祂的诫命，并作为祂司祭的监督，处于某种崇高地位，能禁止非法之事，并教导应遵循之事。因为我们听说，在你们居住的西方地区，魔鬼的诡计如此猛烈地攻击基督的子民，并以如此众多的欺骗爆发出来，不仅压迫和困扰世俗平信徒，也压迫和困扰上主自己的司祭。因此，我们深陷悲痛，不能隐瞒我们应严厉纠正之事。因此，对这些创伤必须采取足够的补救措施，以免治疗上的轻率对头部的致命疾病毫无用处；也免得因过于轻易处理，而因不合法的治疗方式的缺陷，使伤害与治疗者共同陷入邪恶。\par

% structure: p0095 source=h3 target=subsection reason=relative-heading-rank; base=h2

\subsection{一、论那些不应被允许享有控告权的人。}

因此，我们规定并决定，凡品行不端、行为可指责、或信仰、生活及自由状况不明之人，不应有权控告上主的司祭，以免卑鄙之人因此被允许享有控告他们的自由。同样，那些自身涉及任何控告事项或受怀疑之人，也不应有控告长辈的权利；因为受怀疑者和仇敌的声音常会压迫真理。\par

% structure: p0097 source=h3 target=subsection reason=relative-heading-rank; base=h2

\subsection{论外来的裁判。}

此外，通过一项普遍规定，且无损于宗徒在一切事上的权威，我们禁止外来的裁判，因为不应由陌生人审判那本应由本省的人及他自己选定的人作为审判官之人，除非有上诉。因此，若某位主教被明确指控，他应由本省所有主教审理；因为被告不应在本地以外受审。再者，若有人认为他的审判官对他不利，他应主张上诉权；上诉者不应受任何形式的压迫或拘留；但上诉者应有通过上诉补救其受冤案件的权利。在刑事案件中也应有上诉的自由。被判刑者也不应被剥夺上诉权。\par

% structure: p0099 source=h3 target=subsection reason=relative-heading-rank; base=h2

\subsection{三、论被告。}

被告应当在其审判者面前为自己辩护；若审判者非其正当审判者，被告可拒绝发言；每当被告上诉时，应予以宽限（\OriginalTerm{宽限}{induciæ}）。\par

% structure: p0101 source=h3 target=subsection reason=relative-heading-rank; base=h2

\subsection{四、论有人在激动中提出指控或未能证明其所断言之事。}

因此，若有人在激动中轻率地指控某人，则不得将单纯的辱骂视为指控。但应在给予一定时间处理此事后，此人应以书面形式表明其有能力证明自己在激动中所言之事；如此，若他碰巧对自己在激动中的言辞有所反省，并拒绝重复或书写这些言辞，则该人不被视为已被指控犯罪。因此，凡提出指控者，应以书面形式陈述其证明能力。确实，案件应始终在指控被受理之处处理；而无法证实其断言者，应自行承担其提出的刑罚。\par

% structure: p0103 source=h3 target=subsection reason=relative-heading-rank; base=h2

\subsection{五、论被指控的主教向宗座上诉的问题。}

此外，也制定：若被控告的主教向宗座提出上诉，则应以该座教宗的裁定为定案。但在一切有关司铎的案件中，应保持此一形式：任何人不受其本人合法审判官以外之人所宣告的判决之约束。所有信徒也有责任随时准备帮助受压迫和困苦中的人，以便藉着另一种报应方式（\OriginalTerm{报应方式}{vindictæ}）的彰显，能使天主的报应（\OriginalTerm{报应}{vengeance}）远离自己。因为那从受苦者身上除去逆境的人，是向主献上（\OriginalTerm{献上}{libat}）了顺利之事。为此，经上记载：\BibleQuote{兄弟辅助兄弟，必将被高举。} \BibleRef{箴 18:19} 因为天主的教会应当没有瑕疵或皱纹，因此不应受某些人的践踏和玷污；因为经上记载：\BibleQuote{我的鸽子，我的完全人，只有这一个。} \BibleRef{歌 6:9} 此外，主又对梅瑟说：\BibleQuote{在我这里有地方（\OriginalTerm{在我这里}{penes me}），你要站在磐石上。} \BibleRef{出 33:21} 有什么地方不属于主呢？万有既由祂所造，又都在祂内存在。然而，在天主那里确实有一个地方——即圣教会的合一——在那里，人站在磐石上，同时谦卑地持守着\OriginalTerm{宣认的坚固}{confessionis soliditas}。我们劝告你，我们的兄弟，以及我们所有在基督教会中作治理者的兄弟（这教会是祂用自己的血所赎买的），要尽你们所能用一切约束，阻止所有人陷入那深渊，有些弟兄正滑入其中，他们辱骂主的牧者，并用言语和行为迫害他们；我们劝你们不要让他们被情欲的钩子所伤：因为经上记载：\BibleQuote{人的愤怒并不成就天主的正义。} \BibleRef{雅 1:20} 为此又说：\BibleQuote{各人要敏于听，缓于说，缓于怒。} \BibleRef{雅 1:19} 我不怀疑你们依靠天主的帮助遵守这一切；但既然有了劝诫的机会，我也悄悄把我的话加在你们美好的意愿和行动上，好使你们出于自己、未经劝告所做的，如今因有劝告者加入，就不再是独自完成。因此，弟兄们，你们和所有信徒理当彼此相爱，不毁谤或控告对方：因为经上记载：\BibleQuote{你要爱你的邻人，并对他忠信。但你若泄露他的秘密，就不可再追随他。因为人如同毁坏朋友一样，使邻人的爱丧失。又如人从手中放走一只鸟，你放走了邻人，就不得再得回他。不可再追随他，因为他已远去；他如同从网罗逃脱的麋鹿，因为他的灵魂受了伤。你再不能捆住他；辱骂之后或许有和好，但泄露朋友的秘密是不幸灵魂的绝望。以眼示意的人行恶，众人必弃绝他。你在时，他说话甘甜，赞赏你的言语；但最终他会扭动嘴巴，诽谤你的话。我憎恶许多事，但无过于此人；主也必憎恶他。谁向高处投石，石头必落在自己头上；诡诈的打击必使诡诈者受伤。谁挖坑，必掉在其中；谁在邻人路上放绊脚石，必跌在其上；谁为别人设陷阱，必自陷其中。谁作恶，恶必临到他身上，他不知道灾祸从何而来。讥诮和侮辱来自骄傲者；报应如狮子埋伏等候他们。凡因义人跌倒而欢喜的，必被网罗缠住；痛苦在他们死前必吞噬他们。恼怒和愤恨两者都是可憎恶的，罪人必两者兼有。} \BibleRef{德 27:17} \Quote{凡要报仇的，必从主那里得到报应，祂必\Emph{牢记}他的罪恶。你要饶恕邻人伤害你的事，这样你祈祷时，你的罪也必得赦免。人若对人怀恨，岂能从主求宽恕？他对与自己相似的人不施仁慈，岂能向至高者求赦免自己的罪？他既是血肉之人，却怀恨在心，岂能向天主恳求仁慈？谁能为他的罪求饶恕？你要记念你的结局，止息仇恨。因为败坏和死亡威胁着祂的诫命。你要记念对天主的敬畏，不可向邻人怀怨。记念至高者的盟约，宽恕邻人的无知。远避争讼，你便减少罪恶。因为暴躁的人必挑起争端，罪人会扰动朋友，在和平者中制造纷争。正如树林中的树木使火焰燃烧，人的怒气也如此；按人的力量，他的怒气也如此；按他的财富，他的愤怒也升高。仓促的争论点燃烈火，仓促的争斗流血；\OriginalTerm{作证的口舌}{testificans}导致死亡。你若吹火星，它必如火烧；你若吐口水，它必熄灭；两者都从你口而出。暗中诽谤和一口两舌的人被咒骂，因为他毁灭了许多和睦的人。\OriginalTerm{第三者的舌头}{tertia}扰乱了多人，使他们从一国流亡到另一国。它拆毁了富人的坚固城邑，推翻了伟人的房屋；它摧毁了民族的勇力，驱散了强盛的国家。\OriginalTerm{第三者的舌头}{tertia}赶出了\OriginalTerm{贤淑的妇女}{viratas, spirited}，剥夺了她们的劳动成果。谁听从它，永不得安息，永不能有可依赖的朋友。鞭痕留下印记，但舌头的伤痕能折断骨头。许多人死于剑锋，但死于舌头的人更多。那能防御恶舌的人有福了，未经历其毒害；未负过它的轭，未被它的绳索捆绑。因为它的轭是铁轭，它的绳索是铜索。它的死亡是恶性的死亡，坟墓尚且好过它。它的持久不能长久，但它必占有不义之人的道路。在它的火焰中，它不焚烧义人。凡离弃上主的人必陷在其中；它必在他们里面燃烧，永不熄灭；它必如狮子般被派到他们身上，又如豹子吞噬他们。你要用荆棘\OriginalTerm{围住耳朵}{sæpi aures}，拒绝听恶舌；为你的嘴做门，为你的耳朵做闩。你要\OriginalTerm{熔炼}{confla}你的金银，为你的言语做天平，为你的嘴做正直的嚼环。你要谨慎，免得你在舌头上失足，倒在你窥伺的敌人面前，你的跌倒不可救治，直至死亡。} \EditorialNote{德训篇第二十八章} 众人当提防这些事，并且\Quote{要禁止舌头出恶言，嘴唇不说诡诈的话。} \BibleQuote{最后，亲爱的诸位，你们要因主并藉祂大能的力，成为坚强的。穿上天主的全副武装，为能抵抗魔鬼的阴谋，因为我们战斗不是对抗血肉，而是对抗率领者，对抗掌权者，对抗这黑暗世界的统治者，对抗\OriginalTerm{天上}{cœlestibus}邪恶的神灵。为此，你们要拿起天主的全副武装，好在邪恶的日子能够抵挡，并\OriginalTerm{在一切事上完全}{omnibus perfecti}。所以要站稳，用真理束腰，穿上正义的甲，脚上穿着和平福音的预备鞋；在一切事上（\OriginalTerm{在一切事上}{in omnibus}）拿起信德的盾牌，用它就能熄灭恶者所有的火箭。并戴上救恩的头盔，拿着圣神的剑，就是天主的圣言。} \BibleRef{弗 6:10} 兄弟，我们盼望，我们所写给你们的一切，应普遍让众人知晓，为使涉及他人的事也传到众人耳中。愿全能的天主保护你，兄弟，以及我们所有各处的主内弟兄，直到终结——即那位愿意救赎全世界的、我们的主耶稣基督，祂永受赞美。阿们。—— 于十月十六日，在最著名的阿弗里卡努斯与德基乌斯执政期间颁赐。\par

% structure: p0105 source=h2 target=section reason=relative-heading-rank; base=h2

\section{教宗法比盎的谕令}

\Emph{取自 \WorkTitle{格拉齐亚诺谕令集}。}\par

% structure: p0107 source=h3 target=subsection reason=relative-heading-rank; base=h2

\subsection{一、拒绝与兄弟和好者应以最严厉的斋戒来约束。}

如果任何受伤害者拒绝与他的兄弟和好，当伤害他的人做出补偿时，他应以最严厉的斋戒来约束，甚至直到他怀着感恩之心接受所提出的补偿为止。\par

% structure: p0109 source=h3 target=subsection reason=relative-heading-rank; base=h2

\subsection{二、明知而故意发虚誓者被视为可耻。}

凡明知而故意发虚誓者，应处以四十天面包和水的斋戒，并在接下来的七年里做补赎，且永远不可脱离补赎；他永远不得被允许作证。但此后，他可以领受共融。\par

% structure: p0111 source=h3 target=subsection reason=relative-heading-rank; base=h2

\subsection{三、患疯癫的男女不能缔结婚姻。}

患疯癫的男子或女子都不能进入婚姻关系。但如果已经缔结，则不应分离。\par

% structure: p0113 source=h3 target=subsection reason=relative-heading-rank; base=h2

\subsection{四、第五亲等的婚姻关系可以彼此结合；在第四亲等，若发现，不应分离。}

关于通过夫妻关系进入姻亲关系的亲属，在妻子或丈夫去世后，可以在第五亲等缔结结合；在第四亲等，若发现，不应分离。然而，在第三亲等内，一个人在他去世后娶另一个人的妻子是不合法的。同样，男人在其妻子去世后，可以与自己的女亲属以及妻子的女亲属缔结婚姻。\par

\Emph{针对紧接上文的通知。}\par

那些与血亲结婚并被分离的人，只要双方都在世，就不得自由地与其他妻子缔结婚姻，除非他们能以无知作为借口。\par

% structure: p0117 source=h3 target=subsection reason=relative-heading-rank; base=h2

\subsection{五、唯有血亲，或者，若后代完全断绝，年长且可信赖的人，应在主教会议中计算亲属关系事宜。}

任何外人不得控告血亲，或在主教会议中计算血亲关系，而是应由与知晓此事相关的亲属——即父亲、母亲、姊妹、兄弟、伯叔、舅父、姑母、姨母及其子女。然而，若后代完全断绝，主教应按教规向知晓同一关系的年长且可信赖的人询问；若发现此种关系，双方应被分离。\par

% structure: p0119 source=h3 target=subsection reason=relative-heading-rank; base=h2

\subsection{六、每位信徒每年应领圣体三次。}

即使他们不更频繁地领受，但平信徒每年至少应在三次领受圣体，除非某人因更严重的罪行受阻——即复活节、五旬节和主的圣诞节。\par

% structure: p0121 source=h3 target=subsection reason=relative-heading-rank; base=h2

\subsection{七、司铎不应在三十岁以下被祝圣。}

若某人未满三十岁，绝不应被祝圣为司铎，即使他极其堪当；因为主自己也是在三十岁时才受洗，并在此时期开始教导。因此，一位将要被祝圣的人，直到达到此合法年龄之前，不应被祝圣。\par

% structure: p0123 source=h2 target=section reason=relative-heading-rank; base=h2

\section{教宗法比盎的谕令}

\Emph{同一教宗的谕令，取自 \WorkTitle{十六卷谕令法典}，第五卷，第七及第九章。}\par

% structure: p0125 source=h3 target=subsection reason=relative-heading-rank; base=h2

\subsection{一、祭台的奉献应在每个主日进行。}

我们规定，在每个主日，应由男女以饼和酒进行祭台的奉献，以便藉着这些祭献，他们能摆脱自己罪过的重担。\par

% structure: p0127 source=h3 target=subsection reason=relative-heading-rank; base=h2

\subsection{二、不识字（无学识）的司铎不得冒险举行弥撒。}

祭献不得从不胜任按照宗教礼仪履行弥撒中的祷文或行动（\OriginalTerm{行动}{actiones}）及其他礼规的司铎手中接受。\par

\SourceRecord{Translated by S.D.F. Salmond.；Ante-Nicene Fathers；Vol. 8.；Edited by Alexander Roberts, James Donaldson, and A. Cleveland Coxe.；Buffalo, NY: Christian Literature Publishing Co.,；1886.；<http://www.newadvent.org/fathers/0835.htm>.}
